% bitnet-baremetal-t5b.tex
%
% Sub-Two-Bit Ternary Weight Storage with Table-Free Arithmetic Decoding
% Project: bitnet-baremetal   Format: t5b (five trits per byte)
%
% Build:  pdflatex bitnet-baremetal-t5b.tex   (twice, for references)
%
% MECHANISM, IN ONE SENTENCE, because it decides the related work and the
% arXiv category: this is a LOSSLESS RE-ENCODING of already-ternary weights --
% a radix-3 source code at 1.600 bits/weight instead of 2.000 -- together with
% a SIMD kernel that decodes it arithmetically and contracts on the digits, so
% the model is bit-for-bit unchanged and only storage and kernel differ. It is
% not pruning, not weight sharing, and not quantisation.

\documentclass[11pt,a4paper]{article}

\usepackage[utf8]{inputenc}
\usepackage[T1]{fontenc}
% Latin Modern is cosmetic. arXiv's TeX Live has it; a minimal local
% installation may not, and a preprint that will not compile on the machine
% that wrote it is worse than one set in Computer Modern.
\IfFileExists{lmodern.sty}{\usepackage{lmodern}}{}
\usepackage[margin=1in]{geometry}
\usepackage{amsmath,amssymb,amsthm}
\usepackage{booktabs}
\usepackage{array}
\usepackage{listings}
\usepackage{xcolor}
\usepackage{microtype}
\usepackage[hidelinks]{hyperref}
\usepackage{url}
\usepackage{siunitx}
% group-digits=integer is load-bearing: without it siunitx groups the
% FRACTIONAL part too and typesets 1.6875 as "1.687,5", which a European
% reader parses as 1687.5. A reviewer caught it in the compiled PDF.
\sisetup{group-digits=integer,group-separator={,},group-minimum-digits=4}

\newtheorem{proposition}{Proposition}

\lstset{
  basicstyle=\ttfamily\small,
  columns=fullflexible,
  breaklines=true,
  frame=single,
  framesep=4pt,
  backgroundcolor=\color{gray!8},
}

\newcommand{\code}[1]{\texttt{#1}}

\title{Sub-Two-Bit Ternary Weight Storage with\\
       Table-Free Arithmetic Decoding for CPU Inference}
\author{Justus Theile\\
  \small Independent researcher\\
  \small \texttt{justus.theile@gmail.com}\\
  \small \url{https://github.com/purpleskulll/bitnet-t5b}}
\date{September 2026}

\begin{document}
\maketitle

\begin{abstract}
Ternary (1.58-bit) large language models are stored in practice at two bits per
weight, \SI{26}{\percent} above the information content of a ternary symbol.
\textbf{The code that removes that redundancy is not new, and neither is the way
to decode it.} Five ternary values as the base-3 digits of one byte at
\num{1.600} bits per weight is \code{llama.cpp}'s \code{TQ1\_0}, is published
with a correctness theorem as Spectra~1.1's \code{TQ1}~\cite{spectra}, appears a
year earlier in statistical genomics~\cite{miraculix} and fifteen years earlier
in heuristic search~\cite{breyerkorf}; the $32k+j$ interleave that lets a digit
plane address 32 consecutive activations is \code{i2\_s}'s and \code{TQ1\_0}'s
alike; recovering all five digits
\emph{independently} from the packed byte, table-free, is NTRU's since
2019~\cite{ntru} and is what \code{TQ1\_0}'s own NEON path does; contracting on
the digit planes without materialising a weight is \code{TQ1\_0}'s on both
instruction sets; carrying the scale per tensor rather than per block was
\code{block\_q1\_3}, on the pull request that became \code{TQ1\_0}; and the
condition under which a denser format pays is Equation~1 of Zukowski et
al.~\cite{zukowski}, twenty years old. \S\ref{sec:priorart} states all six
concessions and \S\ref{sec:contrib} what remains beside them.

What this paper offers instead is measurement. We build \textbf{t5b} --- the
same construction at a block length of $160=32\times5$, which takes no remainder
byte, with the scale in \code{i2\_s}'s existing 32-byte tail so that no graph
node changes --- and integrate it into \code{llama.cpp} as a distinct
\code{ggml} type alongside the existing two-bit format, then measure what it is
worth. The AVX2 realisation of the depth-one decode through the high half of a
16-bit multiply is the one instantiation upstream does not have, and names as an
open question in its own source. We answer that question by transplanting the
form into upstream's kernel and measuring it there, where it turns out
\emph{not} to pay --- for a reason that is a property of upstream's encoding
rather than of the idea, and that does not transfer to this one. The profitability condition is measured rather than
assumed, one level down the memory hierarchy from where it was first stated,
with both of its terms taken from the part the kernel runs on. So measured, t5b
reduces the shipped \code{BitNet-b1.58-2B-4T} checkpoint
from \SI{1.10}{\gibi\byte} to \SI{1.00}{\gibi\byte} and yields, under
\code{llama-bench} \emph{at four threads}, a median $1.119\times$ prompt and
$1.133\times$ generation throughput over five invocations, faster in five of five. The re-encoding is lossless and
the output is bit-identical: perplexity over \num{20480} tokens of WikiText-2
agrees with the two-bit baseline to six significant figures. We further report
that \SI{4.76}{\percent} of the model's feed-forward neurons are identically
zero, that the dead index sets of the gate and up projections coincide in all
thirty layers, and that an independent fine-tune revives none of them. Separately,
and as a property of the loader rather than of either format, \code{llama.cpp}
builds this architecture's feed-forward block with \code{SiLU} where the
checkpoint declares squared ReLU: supplying the declared activation moves
perplexity over the same tokens from \num{96.8243} to \num{16.7482}, a factor of
\num{5.78}, and the bit-identity above holds under both. That it is the loader
and not the encoding is what the reproduction shows: on upstream
\code{llama.cpp} at \code{TQ1\_0} the same pair reads \num{96.6913} against
\num{16.7750}, so the effect holds to better than \SI{0.2}{\percent} across
three quantisations and two trees whose only shared component is the graph. All
results are from a single AVX2 microarchitecture \emph{without} VNNI. Two
independent runs bound that limitation: a static pipeline model, calibrated
against the host to \SI{3.0}{\percent}, puts the arithmetic ratio at or below
its Zen 2 value on every newer target it covers --- four distinct scheduler
models rather than six, since \code{znver3}, \code{znver4} and \code{znver5}
return identical counts, so what is supported is one-sided, that the ratio does
not run away, and not that it is constant --- and a reviewer's Intel Xeon
reproduces it at \num{2.11} against our \num{2.43}. On real VNNI hardware the same reviewer
measures the ratio rising by \SI{8.8}{\percent} --- the direction predicted, from
the predicted mechanism, at a magnitude smaller than the eight-bit storage choice
of a system built for those targets would suggest, though on a shared host whose
spread exceeds the effect.
\end{abstract}

\section{Introduction}

Quantisation to ternary weights $w\in\{-1,0,+1\}$ preserves task performance at
the two-billion-parameter scale while removing multiplication from the dominant
matrix operation. The resulting models are attractive for CPU deployment: the
matrix--vector product requires only sign application and integer accumulation,
and the weight matrix is small enough to change which hardware resource binds
the computation.

That last point is usually made as a footprint argument and left there. We take
it as the design constraint. During autoregressive generation every weight is
read exactly once per emitted token and used exactly once, so arithmetic
intensity is fixed at one multiply--accumulate per stored weight and the
achievable token time is bounded below by
\begin{equation}
t_{\text{token}} \;\ge\; \max\!\left(\frac{N\mu}{\pi},\;\frac{N\beta}{8B}\right),
\label{eq:roofline}
\end{equation}
where $N$ is the number of weights, $\beta$ the bits used to store each, $\mu$
the instructions issued on the binding execution port per multiply--accumulate,
$\pi$ that port's throughput, and $B$ the memory bandwidth available to the
executing core. The first term is arithmetic; the second is transport.

Existing ternary CPU formats occupy the extremes of the $(\beta,\mu)$ trade-off
that~\eqref{eq:roofline} describes. \code{bitnet.cpp}'s \code{i2\_s} --- the
baseline throughout, and a type of its \code{llama.cpp} fork rather than of
\code{ggml} --- uses
$\beta=2$ with a decode of three shifts and four masks per 128 weights.
\code{bitnet.cpp}'s \code{TL2} and the T-MAC system reach
$\beta\approx\num{1.667}$ by table lookup, paying memory for the tables
(\code{TL1} is at \num{2.0}); \code{llama.cpp}'s own \code{TQ1\_0} reaches
\num{1.6875} by base-3 packing without a table, and \code{TQ2\_0}
\num{2.0625}; Spectra~1.1~\cite{spectra} publishes that family with a
correctness theorem and CPU kernels as \code{TQ1} and \code{TQ2}, recommending
exactly the $p=8$, $k=5$ this work uses. Every format named so far is lossless.
In the lossy direction, \emph{Sherry}~\cite{sherry} reaches \num{1.25} bits by
imposing 3:4 fine-grained sparsity on the ternary weights and packing four of
them into five bits, restoring power-of-two alignment at the cost of a
training-time constraint and of exactness --- a different trade from the one
studied here, where the weights are unchanged bit for bit. In the
opposite direction, \emph{Litespark} stores ternary weights at $\beta=8$, on the
explicit grounds that ``using 2-bit packing would require unpacking weights
before every computation, negating the performance benefit'' --- a decision made
for targets whose dot-product instruction (\code{VPDPBUSD}, ARM \code{SDOT})
consumes eight-bit operands.

\paragraph{The base-3 packing itself is not new.} \code{llama.cpp} has carried \code{TQ1\_0} since August
2024 (pull request~\#8151): its block structure in \code{ggml-common.h} is
annotated \emph{``5 elements per byte ($3^5 = 243 < 256$)''}, it decodes without
a table, and it has an AVX2 dot product. \code{TQ1\_0} reaches \num{1.6875}
bits per weight only because each 256-element block carries an f16 scale and a
16-element remainder at two bits; the code bytes themselves are at \num{1.600}.
Any claim of novelty for the packing would survive on the arithmetic of that
block overhead alone, which is not a claim worth making. We had integrated into
this very codebase without checking it --- and the same holds for the extraction
arithmetic and for the scale convention conceded in \S\ref{sec:priorart}. All
three were found after the fact. This paper records that here, once, rather than
in a footnote.

What this paper contributes is therefore narrower than any of that, and
\S\ref{sec:priorart} concedes each piece to its owner. What is left: a block
length of $160=32\times5$, which takes no remainder byte where \code{TQ1\_0}'s
\num{64} does; a scale carried in \code{i2\_s}'s own 32-byte tail, so the
format is drop-in for the \code{i2\_s} path of \code{bitnet.cpp} with no graph
node and no loader change; the AVX2 instantiation of a depth-one decode that
upstream has only on NEON; a column-blocked matrix--matrix kernel, which the
nearest prior format has no equivalent of; the \emph{measurement} of a
twenty-year-old condition at a level of the hierarchy where it had not been
measured; and an end-to-end evaluation on a real ternary checkpoint against the
format that ships with it.

\paragraph{What this is and is not.}
The mechanism is a \emph{lossless re-encoding} of weights that are already
ternary, together with a decode-fused kernel. No weight value changes; the
output is bit-identical. It is therefore neither pruning nor weight sharing nor
a new quantisation scheme, and it composes with all three.

It is also not a cache format, and that boundary is where the largest remaining
memory sits. \S\ref{sec:ratios} measures the part of a token this work does not
touch at \SI{21.6}{\milli\second}, the larger half, and the key/value cache
lives entirely inside it. Hariri et al.~\cite{spectralkv} act on that half rather
than this one: they derive from the spectral and Frobenius norms of the key and
value projections that keys deserve more bits than values, and allocate four and
two, keeping up to \SI{98.3}{\percent} of a uniform allocation's accuracy at
lower memory. The two address disjoint tensors --- weights here, activations
there --- so they compose without interacting, which is a statement about scope
and not a measured result; we have not run the combination.

\subsection{Prior art, and what is not claimed here}\label{sec:priorart}

Six constructions this work is built from belong to someone else. They are
conceded here, once and in one place: the base-3 packing, the $32k+j$
interleave, the depth-one extraction arithmetic, the fused digit-plane
contraction, the per-tensor scale model and the profitability condition.

\paragraph{The packing.} Five ternary values as the base-3 digits of a byte
is not ours: it is \code{TQ1\_0}'s construction in \code{llama.cpp}, and
it is published with a correctness theorem by Vaidhya et
al.~\cite{spectra} as Spectra~1.1's \code{TQ1}, whose $p=8$, $k=5$
recommendation and \num{1.6}-bit figure are the same as ours. Nor does it
originate in this field: Breyer and Korf store heuristic estimates modulo
three at \num{1.6} bits per entry for pattern databases in heuristic
search~\cite{breyerkorf}, fifteen years before ternary language models,
and the \code{5codes} format of \emph{miraculix} publishes the same
construction in statistical genomics a year before
\code{TQ1\_0}~\cite{miraculix} --- its Algorithm~1 writes out the
positional weights $3^{0},\dots,3^{4}$, its source defines
\code{BitsPerCode 1.6}, and its stated motivation is the one taken here,
to keep the multiplication on the compressed data so that nothing is
decompressed. Both decode by table: miraculix precomputes the $3^{5}$
possible partial dot products per tile and uses the packed byte directly
as the index, so it belongs with the table-lookup systems above rather
than with the decode of \S\ref{sec:kernel}.

\paragraph{The interleave.} The $32k+j$ layout that makes digit plane $k$
address 32 consecutive activations is \code{i2\_s}'s at two bits and
\code{TQ1\_0}'s at five digits per byte, on the same instruction set.
\S\ref{sec:format} states it against both, because it is what the drop-in claim
rests on.

\paragraph{The scale model.} \code{block\_q1\_3} lived on \#8151 from 2024-06-19 to
2024-07-30: five elements per byte, annotated in the same words
(\emph{``5 elements per byte ($3^5 = 243 < 256$)''}) and carrying
\emph{no} scale field at all, at \num{1.625} bits per weight. Its scale
was a separate per-layer \code{ggml} tensor of extent one applied after
the matmul --- \code{Qcur = ggml\_mul(ctx0, Qcur,
model.layers[il].wq\_scale)} in \code{build\_bitnet}. It was replaced by
\code{TQ1\_0}'s per-block f16 deliberately, so that the stock computation
graph could be reused, and the cost of that choice was recorded at the
time as wasting \num{0.0625} bits per weight. The same construction ships
today as \code{IQ1\_BN} in \code{ik\_llama.cpp}.

What is left of the format is narrower again, and worth stating with its true
size. Of the \num{20828160} bytes t5b saves against \code{TQ1\_0},
\code{Q1\_3} already delivered \SI{78.2}{\percent}; the residual is
\num{4546560} bytes, \SI{1.03}{\percent} of the \code{TQ1\_0} ternary
payload and \SI{0.38}{\percent} of the file. It is a block-length choice
rather than a scale model: $64 = 12\times5+4$ needs a remainder byte and
lands at \num{1.6250}, whereas $160 = 32\times5$ needs none and lands at
\num{1.6000}. The one thing here that is not anticipated is that the
scale rides in \code{i2\_s}'s own 32-byte tail, so no graph node and no
loader change is required where \code{Q1\_3} needed both --- and that is
\code{i2\_s}'s convention, not ours.

\paragraph{The extraction arithmetic.} That a
multiplication-based decode avoids division and modulo is also known:
Spectra~1.1 exploits $3^{5}\approx 2^{8}$ to extract the trits
\emph{iteratively}, $b_{i+1}=(3b_{i})\wedge\texttt{0xFF}$ for
$i=0\dots4$, a strictly serial chain of five dependent multiplies. The
decode used here instead takes all four quotients
$\lfloor x/3^{j}\rfloor$ independently, straight from the original
byte, at one \code{vpmulhuw} each, which is exact because a five-trit
byte satisfies $x\le\num{242}$ and so lies inside the exactness range of
16-bit magic multipliers --- collapsing the depth-five dependence to depth
one. That form is not ours either, and at exactly these
parameters. The NTRU reference implementation has recovered all five
trits of a packed byte since 2019 with four independent multiplies read
directly from that byte --- constants \num{171}, \num{57}, \num{19},
\num{203} with shifts \num{9}, \num{9}, \num{9}, \num{14}, which we
verified to be exact floor divisions by $3,9,27,81$ over all 256 byte
values --- at dependency depth one and without a lookup
table~\cite{ntru}.\footnote{\code{ref-common/pack3.c}, function
\code{poly\_S3\_frombytes}. It does contain a serial divide-by-three
chain, but only in a leftover tail of at most three coefficients, which is
not compiled at all for the \code{hps4096821} and \code{hrss701} parameter
sets.}

\paragraph{The fused contraction.} NTRU materialises every coefficient into a
coefficient array
and defers the modulo-three reduction to a separate branch-free sweep,
where the kernel here feeds the digit planes straight into the
contraction. That distinction does not rescue a novelty claim
either, and the counter-example is the kernel this paper positions itself
against. \code{TQ1\_0}'s NEON path derives all four scalings
$3,9,27,81$ of the loaded byte with one \code{vmulq\_u8} each, every one
of them from the \emph{original} vector rather than from its predecessor,
extracts the digits table-free by multiplying by three and keeping the two
bits above eight, and passes them directly to \code{vdotq\_s32}. That is
depth one, table-free, and without materialisation, and it has shipped
since pull request~\#8151. Its AVX2 path does the same without
materialisation but at depth \emph{two}: lacking an 8-bit multiply it
chains the scalings as shift-and-add, deriving $27$ from $3$ and $81$ from
$9$ --- except in its four-element tail, which reaches depth one through
\code{\_mm256\_mullo\_epi16} against a vector of $1,3,9,27$.

What is therefore left is one instantiation, not a mechanism: the depth-one
form in the AVX2 \emph{main} path, reached through the high half of a
16-bit multiply. Upstream states that as an open question in the file
itself --- \code{// TODO: can \_mm256\_mulhi\_epu16 be faster even if
16-bits?}, \code{arch/x86/quants.c:1406}. We verified every sentence of
this paragraph against the vendored kernel at \code{390c307} rather than
against a description of it. We answered that question by transplanting the
form into upstream's kernel, and the answer is no --- for \code{TQ1\_0};
\S\ref{sec:tq10} reports the measurement and why its reason does not transfer
to this encoding.

Against the serial form \emph{this} encoding admits, depth
one is free rather than expensive. A chain
$q_{j+1}=\code{mulhi}(q_{j},\num{21846})$ over the two views is two
multiplies at each of four levels, the same eight the independent form
issues; what changes is only the dependence. Nor does the latency it buys
bind: \S\ref{sec:results} measures the loop at \SI{67}{\percent} of its
multiply-port ceiling. So the honest statement is that depth one costs
nothing and relieves a constraint that was not the binding one, which is a
weaker claim than an earlier draft made. No head-to-head measurement
against Spectra's kernel was made: their figures are from a Mac~M4 and an
EPYC part, ours from one Zen~2 host.

\paragraph{The condition.} The criterion is twenty years old: Equation~1 of
Zukowski et al.~\cite{zukowski} switches between the I/O-bound and CPU-bound
regimes on a test that is $S\le\Sigma$ rearranged, and applies it as an
adopt-or-reject rule for a codec.

\subsection{Contributions}\label{sec:contrib}
Each item names the construction it builds on; \S\ref{sec:priorart} concedes
that construction to its owner.
\begin{enumerate}
\item A \num{1.600}-bit variant of the known base-3 packing
      (\S\ref{sec:format}), on \code{TQ1\_0}'s construction and
      \code{block\_q1\_3}'s scale model: block length $160=32\times5$, which
      takes no remainder byte, with the scale in \code{i2\_s}'s own 32-byte
      tail so that no graph node and no loader change is needed.
\item The AVX2 instantiation of a depth-one decode (\S\ref{sec:kernel}), on
      NTRU's extraction arithmetic and \code{TQ1\_0}'s digit-plane
      contraction: the form upstream has on NEON and names as an open question
      in its own AVX2 source, reached here through the high half of a 16-bit
      multiply, with an exactness proposition and an exhaustive
      first-failure table behind it. \S\ref{sec:tq10} answers upstream's
      question by transplanting the form into its kernel, where it does not pay.
\item A column-blocked matrix--matrix kernel (\S\ref{sec:kernel}), which the
      nearest prior format has no equivalent of and which is what carries the
      prompt-processing result.
\item A \emph{measurement} of the profitability condition
      (\S\ref{sec:condition}), on Zukowski et al.'s criterion, derived
      from~\eqref{eq:roofline} and verified by
      measuring both $\pi$ and $B$ rather than inferring either. What is new is
      the level of the
      hierarchy at which it is measured --- there $B$ is disk bandwidth against
      a RAID array and main memory is named only as future work, here it is
      DRAM-to-SIMD --- and the fact that both of its terms are measured on the
      part in question rather than assumed.
\item Integration and end-to-end evaluation (\S\ref{sec:results}) in
      \code{llama.cpp}, with the two-bit path retained as a control and
      losslessness verified by perplexity agreement over
      $\num{2.05}\times10^{4}$ tokens.
\item A structural observation (\S\ref{sec:dead}): \SI{4.76}{\percent} of the
      checkpoint's feed-forward neurons are identically zero, the dead index
      sets of $W_{\text{gate}}$ and $W_{\text{up}}$ coincide in all thirty
      layers, and an independent fine-tune revives none of them.
\end{enumerate}

\section{Background}

\subsection{The information bound}
A ternary weight carries at most
\begin{equation}
H_{\max}=\log_{2}3=\num{1.58496}\ \text{bits},
\end{equation}
the origin of the ``1.58-bit'' label, attained only for a uniform distribution.
The empirical marginal over all \num{2084044800} ternary weights of the shipped
checkpoint is
\begin{equation}
(p_{-1},p_{0},p_{+1})=(\num{0.288940},\ \num{0.422109},\ \num{0.288951}),
\end{equation}
giving a zeroth-order entropy
\begin{equation}
H=-\sum_{s}p_{s}\log_{2}p_{s}=\num{1.5603}\ \text{bits}.
\label{eq:entropy}
\end{equation}
The excess zero mass \emph{lowers} the floor below $\log_{2}3$. Storage at
$\beta=2$ exceeds~\eqref{eq:entropy} by \SI{28.2}{\percent}; the code of
\S\ref{sec:format} exceeds it by \SI{2.5}{\percent}.

That residual \SI{2.5}{\percent} is the one place an entropy coder could still
buy something, and the closest work is Tan et al.~\cite{shannonllm}, which pairs
an entropy study of weights from 1.5\,B to 405\,B parameters with a tile-level
Asymmetric Numeral Systems decoder fused into the GEMM pipeline of GPU
inference, reporting bit rates within \numrange{0.01}{0.1} bits of the bound.
It does not settle the question here, and the reason is worth being precise
about: their format space bottoms out at INT4 and FP4 and treats no ternary
format, and their central observation --- that nominal formats over-allocate by
a factor of 2 to 10 --- does not carry over to a code already within
\SI{2.5}{\percent} of its own measured entropy. The obstacle on this target is
not the coder's rate but its cost: a variable-rate decoder must still satisfy
\eqref{eq:condition} against a kernel sustaining \num{632} million
128-weight blocks per second, and their decoder
reconstructs tiles into shared memory for tensor cores rather than contracting
without materialising. Whether an interleaved rANS decoder can meet that bound
on a SIMD CPU is open, and this paper does not answer it.

\subsection{The interface a replacement must preserve}
\code{ggml}'s ternary kernels do not compute $\sum_i w_ia_i$ directly. They
compute on the shifted code $c_i=w_i+1\in\{0,1,2\}$, the unsigned operand
\code{vpmaddubsw} requires, and the caller removes the shift:
\begin{equation}
\sum_{i=1}^{n}w_ia_i=\sum_{i=1}^{n}c_ia_i-\sum_{i=1}^{n}a_i .
\label{eq:shift}
\end{equation}
The correction depends only on the activations and is therefore evaluated once
per matrix--vector product rather than once per row, $O(n)$ against $O(mn)$. Any
replacement must expose the same quantity $\sum_i c_i a_i$. Identity
\eqref{eq:shift} requires $c_i\le 2$; a code $c=3$ would contribute $3a_i$ while
denoting $w=0$. We verify in \S\ref{sec:correct} that code~3 occurs zero times
across all \num{2084044800} weights.

\section{The t5b code}\label{sec:format}

\subsection{Definition}
A block occupies 32 bytes and carries $L=160$ ternary weights. For byte index
$j\in\{0,\dots,31\}$ and digit index $k\in\{0,\dots,4\}$,
\begin{equation}
b_{j}=\sum_{k=0}^{4}c\!\left(w_{32k+j}\right)3^{k},
\qquad c(w)=w+1\in\{0,1,2\}.
\label{eq:pack}
\end{equation}
The interleave $32k+j$ is not ours at either width. \code{i2\_s} uses it with
four two-bit fields: \code{dequantize\_row\_i2\_s} sends bits $6-2k$ of byte
$j$ to weight $32k+j$, and its AVX2 dot pairs those four shifts with activation
loads at $+0,+32,+64,+96$. \code{TQ1\_0} uses it with \emph{five} base-3
digits, in the same file the open question of \S\ref{sec:priorart} comes from:
\code{quantize\_row\_tq1\_0\_ref} packs \code{x[m + n*32]} under the comment
\emph{``5 elements per byte, along 32 bytes''}, and its AVX2 dot contracts five
digit planes of one 32-byte load against activations at
$+0,+32,+64,+96,+128$ --- the same stride on the same instruction set. What
differs is smaller than it looks and matters for \S\ref{sec:kernel}: both
upstream formats put the earliest weight in the \emph{most} significant field
where \eqref{eq:pack} puts it in the least, and \code{TQ1\_0} stores
$\lceil 256v/243\rceil$ rather than the base-3 value $v$ --- which is why its
digits come out by a different arithmetic, and why the exactness argument of
\S\ref{sec:kernel} applies here and not there. We adopt the interleave rather
than invent it, and state it because it is what the drop-in claim rests on:
digit plane $k$ addresses 32 \emph{consecutive} activations, so the access
pattern is unchanged from \emph{both}. The code is
injective since $3^{5}=243\le 256$, and
\begin{equation}
b_{j}\;\le\;2\sum_{k=0}^{4}3^{k}=2\cdot 121=242 .
\label{eq:bound242}
\end{equation}
Density is
\begin{equation}
\beta=\frac{8\ \text{bits}}{5\ \text{weights}}=\num{1.600},
\end{equation}
a factor $F=2/\num{1.6}=\num{1.25}$ fewer bytes than \code{i2\_s}.

\subsection{Optimality among aligned radix-3 codes}
Packing $k$ trits into $\lceil k\log_{2}3\rceil$ bits gives
$\beta(k)=\lceil k\log_{2}3\rceil/k$. Restricting to codes that fit a machine
word without straddling:

\begin{center}
\begin{tabular}{rrrl}
\toprule
$k$ & bits & $\beta$ & fits \\
\midrule
4  & 7  & 1.750 & byte, 1 bit wasted \\
5  & 8  & \textbf{1.600} & byte, exactly \\
9  & 15 & 1.667 & 16-bit word, 1 bit wasted \\
10 & 16 & \textbf{1.600} & 16-bit word, exactly \\
\bottomrule
\end{tabular}
\end{center}

$k=5$ and $k=10$ tie at \num{1.600} and nothing byte- or word-aligned does
better; $k=10$ is admissible because $3^{10} = \num{59049} \le \num{65536}$. We
implement both; \S\ref{sec:lane} shows they differ by $1.5\times$ in throughput
for reasons of SIMD lane width alone, which is why the byte variant is the one
deployed.

\paragraph{Relation to Spectra 1.1's Theorem 1.} Vaidhya et al.~\cite{spectra}
state the \emph{correctness} condition for this family: packing $k$ trits into a
$p$-bit integer is lossless if and only if $2^{p} > 3^{k}$, and they likewise
recommend $p=8$, $k=5$ for an effective \num{1.6} bits. The table above is the
\emph{optimality} counterpart --- which admissible $(p,k)$ minimises $p/k$
subject to $p$ being a whole machine word --- and every entry in it satisfies
their condition by construction. We claim no priority over the condition, the
recommendation or the resulting bit rate: all three are theirs, and $k=5$ is
\code{TQ1\_0}'s construction in \code{llama.cpp} before either. What this
subsection adds is the alignment argument making $k=10$ the only other candidate
at the same density, which \S\ref{sec:lane} then rules out on lane width.

\subsection{Tail}
For $n$ not a multiple of $L$ the final block's unused slots take $c=1$
($w=0$) and the corresponding activations are read from a zero-padded buffer, so
both terms of~\eqref{eq:shift} are unaffected. Storage is
\begin{equation}
S(n)=32\left\lceil \frac{n}{160}\right\rceil\ \text{bytes},
\end{equation}
exactly \num{1.600} bits/weight at $K=\num{2560}$ and \num{1.6296} at
$K=\num{6912}$; over the whole checkpoint \num{1.6076} including tail waste and
the 32-byte per-tensor scale record \code{ggml} appends.

\section{Kernel}\label{sec:kernel}

\subsection{Digit recovery without a chain}
With $q_{j}=\lfloor b/3^{j}\rfloor$ the digits satisfy
\begin{equation}
d_{j}=q_{j}-3q_{j+1},\qquad j=0,\dots,4,\quad q_{5}=0 .
\label{eq:digits}
\end{equation}
A naive implementation derives $q_{j+1}$ from $q_{j}$, a serial dependence of
length five. We obtain every $q_{j}$ directly from $b$, which
\eqref{eq:bound242} makes possible.

\begin{proposition}[sufficient range for exact magic division]
Let $D\ge 2$ and let $M$ be any positive integer with $e:=MD-2^{16}\ge 1$. Then
\begin{equation}
\left\lfloor\frac{xM}{2^{16}}\right\rfloor=\left\lfloor\frac{x}{D}\right\rfloor
\qquad\text{for every integer } 0\le x<\frac{2^{16}-e(D-1)}{e}.
\label{eq:magic}
\end{equation}
\end{proposition}

\begin{proof}
Write $x=qD+r$ with $0\le r<D$. Since $MD=2^{16}+e$,
\[
xM=q\left(2^{16}+e\right)+rM=q\,2^{16}+\left(qe+rM\right),
\]
so $\lfloor xM/2^{16}\rfloor=q+\lfloor (qe+rM)/2^{16}\rfloor$ and it suffices
that $qe+rM<2^{16}$. Using $q\le x/D$ and $rM\le (D-1)M=2^{16}+e-M$,
\[
qe+rM\;\le\;\frac{xe}{D}+2^{16}+e-M ,
\]
which is below $2^{16}$ exactly when $xe/D<M-e$. Substituting
$M=(2^{16}+e)/D$ gives $xe<2^{16}-e(D-1)$.
\end{proof}

The bound is sufficient, not tight. Since the operative question is finite ---
does the identity hold for every $x\le\num{242}$? --- we also determine the
exact first-failure point by exhaustive search over the whole 16-bit domain, and
it is the exhaustive figure the implementation and its tests pin:

\begin{center}
\begin{tabular}{rrrrrrr}
\toprule
$j$ & $D=3^{j}$ & $M$ & $e$ & bound~\eqref{eq:magic} & first failure & margin over 242 \\
\midrule
1 & 3  & \num{21846} & 2   & $x<\num{32766}$ & $x=\num{32768}$ & $135.4\times$ \\
2 & 9  & \num{7282}  & 2   & $x<\num{32760}$ & $x=\num{32768}$ & $135.4\times$ \\
3 & 27 & \num{2428}  & 20  & $x<\num{3251}$  & $x=\num{3293}$  & $13.6\times$ \\
4 & 81 & \num{811}   & 155 & $x<\num{343}$   & $x=\num{485}$   & $2.00\times$ \\
\bottomrule
\end{tabular}
\end{center}

The tightest case first fails at $x=485$, exactly twice the largest legal packed
byte. All four quotients are therefore exact for every input the code can
produce, each costs one \code{vpmulhuw}, and --- crucially --- the four are
mutually independent, so the depth-five chain collapses to depth one and latency
ceases to bind. What binds instead is the total vector-operation count rather
than the multiply port: \S\ref{sec:results} measures the loop at
\SI{67}{\percent} of its multiply-port ceiling, so that port is idle a third of
the time, and \S\ref{sec:shufdig} finds that removing six vector operations of
which only two are multiply-port buys \num{1.2254}, above what either the
multiply-port ratio (\num{1.18}) or the total-op ratio (\num{1.16}) predicts
alone. Note $M=811$ is not the
canonical $\lceil 2^{16}/81\rceil=810$; the implementation uses the larger
constant, which has the smaller exactness range (first failure at \num{485}
rather than \num{806}) and is still twice what the code can produce. The
proposition covers it because its hypothesis is only $MD>2^{16}$. The suite
verifies all four identities exhaustively and reports each first-failure point,
so the margin is a recorded quantity rather than an assumption.
Digits then follow from~\eqref{eq:digits} using shifts and subtractions
($3q=(q\ll 1)+q$) on the unconstrained integer ports.

\subsection{Contraction and the accumulator bound}\label{sec:accum}
Digit plane $k$ is a register of 32 bytes in $\{0,1,2\}$, exactly the unsigned
operand \code{vpmaddubsw} expects. That range is a consequence of
\eqref{eq:bound242} and not a property the decode enforces: the thirteen byte
values $243\ldots255$, which no packer of this format emits, all yield $d_4=3$,
and a block of them would bound a lane at $2\cdot128\cdot11=\num{2816}$, which
the twelve-block fold derived below does not survive. What the subtraction needs is
weaker and does hold unconditionally --- no digit subtraction borrows for any of
the 256 byte values.
\begin{equation}
\text{acc}_{16}\mathrel{+}=\code{vpmaddubsw}\!\left(d_{k},\,a_{32k..32k+31}\right).
\end{equation}
A product is at most $2\times 128=256$ --- the bound the header and the code
use, covering $a=-128$ --- and \code{vpmaddubsw} sums adjacent
pairs, so a 16-bit lane grows by at most 512 per plane and $5\times 512=2560$
per block. Folding to 32-bit every $\Phi$ blocks is safe iff
\begin{equation}
2560\,\Phi\le\num{32767}\iff\Phi\le\num{12.80},
\label{eq:fold}
\end{equation}
hence $\Phi=12$. This is a \emph{worst-case} bound. The \code{i2\_s} reference
kernel folds every 32 blocks with four planes. Its own operands mask the codes
to $0\ldots3$, so a lane can reach $128\times 762=\num{97536}$; the encoder
emits at most code 2, which still gives $128\times 508=\num{65024}$. Either
figure is twice the \num{32767} ceiling or more, and the kernel is correct only
because signed activations cancel. (The bound for t5b above takes $|a|\le128$, covering $a=-128$; the
figures for the reference kernel take \num{127}, which is the convention its
own comments use. The difference is under \SI{1}{\percent} and neither changes
a conclusion, but the two are not the same convention.)
\S\ref{sec:correct} shows that removing the cancellation makes it
disagree with exact arithmetic in \num{11998} of \num{12000} rows at
$K=\num{6912}$ --- statistical rather than absolute, for the margin computed
there.

\subsection{Why bytes and not words}\label{sec:lane}
The ten-trits-per-\code{uint16} code admits a kernel that never forms a digit.
Substituting~\eqref{eq:digits} into $\sum_k d_ka_k$ and re-indexing,
\begin{equation}
\sum_{k=0}^{K-1}d_{k}a_{k}=q_{0}b_{0}+\sum_{j=1}^{K-1}q_{j}b_{j},
\qquad b_{0}=a_{0},\quad b_{j}=a_{j}-3a_{j-1},
\label{eq:telescope}
\end{equation}
so only the quotient chain and a transformed activation vector are needed, the
latter computed once per matrix--vector product. The identity holds modulo
$2^{32}$ rather than term by term: intermediate products are large and cancel.

It is nonetheless slower by $1.5\times$, for reasons of lane width rather than
radix. On \code{uint16} lanes the contraction must use \code{vpmaddwd}, covering
16 lanes where \code{vpmaddubsw} covers 32, and the quotient chain competes for
the same port. Per 160 weights the word variant issues 19 multiply-port
instructions against the byte variant's 13, measuring \num{30.04} against
\num{33.74}~GMAC/s per core (\S\ref{sec:results}).

\subsection{Blocking for batched operation}
For $n_{c}$ activation columns, decoding per column costs $n_{c}$ decodes.
Blocking so a decoded plane is contracted against $N_{C}$ columns held in
registers reduces the per-block cost from $N_{C}(\delta+5)$ to $\delta+5N_{C}$
multiply-port instructions, $\delta=8$. At $N_{C}=8$ this is 48 instructions per
\num{1280} multiply--accumulates, \num{26.7} MACs each against \code{i2\_s}'s
\num{32}. Measured at $N_{C}\in\{2,4,8\}$: \num{50.51}, \num{63.54},
\num{73.26}~GMAC/s. We deploy $N_{C}=8$; its register pressure (181 stack
references per loop) costs less than the amortisation gains.

\section{When a denser code pays}\label{sec:condition}

Let baseline $A$ and candidate $B$ store the same weights at
$\beta_{A}>\beta_{B}$ with $\mu_{A}<\mu_{B}$ binding-port instructions per
multiply--accumulate, and let $A_{A}>A_{B}$ be the arithmetic rates the two
kernels actually attain with memory taken out of the measurement
(\S\ref{sec:results}). Write
\begin{equation}
F=\frac{\beta_{A}}{\beta_{B}},\qquad S=\frac{A_{A}}{A_{B}},
\end{equation}
and define the \emph{surplus} of the baseline,
\begin{equation}
\Sigma=\frac{\beta_{A}/(8B)}{1/A_{A}}=\frac{A_{A}\beta_{A}}{8B} .
\label{eq:surplus}
\end{equation}
If $\Sigma\le 1$ the baseline is arithmetic-bound and no $S>1$ can win. If
$\Sigma>1$ it is transport-bound and, from~\eqref{eq:roofline}, $B$ is faster iff
\begin{equation}
\boxed{\;S<\Sigma\;}
\label{eq:condition}
\end{equation}
in which case it runs a factor $\min(F,\Sigma/S)$ faster.

\paragraph{$A$ is the measured rate, not the port ceiling, and the difference
decides cells.}
\eqref{eq:roofline} bounds the arithmetic term by $N\mu/\pi$, and it is tempting
to read $A_{A}=\pi/\mu_{A}$ straight off the port count. That is the ceiling, not
the rate, and neither kernel reaches it: \S\ref{sec:results} measures the
baseline at \SI{63}{\percent} of its \num{130.9}\,GMAC/s ceiling and the packed
kernel at \SI{67}{\percent} of its \num{50.5}. Substituting ceilings would give
$S=\mu_{B}/\mu_{A}=\num{2.60}$ against the \num{2.43} measured, and would raise
each $\Sigma$ below by that arm's own shortfall rather than by one common
factor, because the shortfall is not constant in thread count: the baseline
attains \num{83.70}\,GMAC/s per core at one thread and \num{71.89} at four
against the same \num{130.9} ceiling, so those two cells would move to
\num{2.12} and \num{3.39} respectively. That reverses the prediction at two and at four
threads, where a loss is what \S\ref{sec:traffic} measures. Both terms are
therefore taken from measured kernel rates throughout, and $\pi/\mu$ appears
only as the ceiling each kernel fails to reach. The condition is a property of a
kernel pair on a machine rather than of the instruction set alone, and both of
its terms are measurable in minutes.

\paragraph{The criterion is not new; the measurement is.}
Equation~1 of Zukowski et al.~\cite{zukowski} selects between an I/O-bound and a
CPU-bound regime on the test $Br/C + Br/Q \le 1$, which is
\eqref{eq:condition} rearranged, with $S=1+Q/C$ the slowdown decode imposes on
the compute side and $\Sigma=Q/(Br)$ the delivery-to-compute ratio at the
compressed footprint. That work already uses it as an adopt-or-reject rule,
solving for a break-even decompression bandwidth of \SI{883}{\mega\byte\per
\second} and rejecting the codecs below it, and its \code{PFOR} family is
designed for the same reason this decode is --- no branches and no dependence
between the values being decompressed, so that decoding pipelines. We claim no
novelty for the criterion. What differs here is the level of the hierarchy:
their $B$ is disk bandwidth against a RAID array and they name main memory only
as future work, whereas both terms below are measured DRAM-to-SIMD on the part
the kernel runs on. Compression-aware rooflines have since been drawn for
PDE solvers~\cite{goetschel} and, with an area axis rather than a time axis, for
CPUs with in-core matrix engines~\cite{deca}; the thread-dependence of the
crossover --- that a denser format can fail at one core and pay at six --- is
likewise a published result in sparse linear algebra~\cite{aliaga}.

\subsection{Measuring $\pi$}
Eight independent dependence chains, so latency cannot bind, with the clock
derived from a dependent add chain retiring one per cycle by construction:

\begin{center}
\begin{tabular}{lrr}
\toprule
instruction & Gop/s & ops/cycle \\
\midrule
\code{vpmaddubsw} & 4.09 & 1.05 \\
\code{vpmaddwd}   & 4.10 & 1.05 \\
\code{vpmulhuw}   & 4.09 & 1.05 \\
\code{vpmullw}    & 4.08 & 1.04 \\
\code{vpand}      & 15.43 & 3.95 \\
\code{vpaddw}     & 11.72 & 3.00 \\
\bottomrule
\end{tabular}
\end{center}

Measured clock \SI{3.91}{\giga\hertz}. All four multiply-class instructions land
on one port at 1.05 per cycle; the rest are three to four times more plentiful.
This fixes $\pi=\num{4.09e9}\,\mathrm{s^{-1}}$ and identifies the binding
resource unambiguously.

\paragraph{Methodological note.} An earlier version of this measurement gave all
eight chains identical seeds; the compiler proved them equal and kept one,
eliminating \code{vpand} entirely and strength-reducing \code{vpaddw} into a
multiply. It reported $\num{2.98e6}$ and \num{9.15} ops/cycle against a ceiling
near 4. Only the physical impossibility of those figures exposed the defect.

\subsection{Measuring $B$: the surplus is a function of core count}
$B$ is a core's \emph{share}, not the socket's. Running the baseline kernel at
\SI{256}{\kibi\byte} per thread (inside L2) and \SI{64}{\mebi\byte} per thread
(four times one CCX's L3), one thread pinned per physical core:

\begin{center}
\begin{tabular}{rrrr}
\toprule
threads & arithmetic (GMAC/s) & transport (GB/s) & $\Sigma$ \\
\midrule
1 & 83.70  & 15.47 & 1.35 \\
2 & 162.42 & 24.94 & 1.63 \\
3 & 235.14 & 33.00 & 1.78 \\
4 & 287.54 & 38.62 & 1.86 \\
6 & 471.74 & 38.86 & \textbf{3.03} \\
\bottomrule
\end{tabular}
\end{center}

Arithmetic scales $5.6\times$ across six cores, per core falling only
\SI{6}{\percent}. Transport scales $2.5\times$ and is flat from four cores to
six. The surplus therefore grows with core count for reasons independent of any
kernel.

\paragraph{$\Sigma_{6}$ is not a constant, and $S$ sits inside its range.}
The six-thread entry is the only cell in which the condition is satisfied, so it
carries the argument, and the figure above is a single run on a quiet host.
Repeating it five times at load averages of \numrange{3.8}{7.0} gives
\num{1.65}, \num{1.69}, \num{1.75}, \num{1.92} and \num{2.38}, with a median of
\num{1.75}. Those runs do not refute \num{3.03}: under contention a
six-thread arm does not get six cores, and the cache-resident numerator suffers
more from that than the bandwidth-limited denominator, so a contended
measurement is a lower bound. What they establish is that $\Sigma_{6}$ is a
property of the machine's occupancy as well as of the part, and that
$S=\num{2.43}$ sits \emph{above every loaded measurement and below the idle
one}. At \num{3.03} the denser format wins at six threads; at any of the five
contended values it does not. The sign of the
one positive prediction therefore depends on the host being idle, which is a
weaker statement than the table alone suggests
(\code{results/sigma\_six\_threads.txt}).

\paragraph{What \eqref{eq:condition} therefore predicts, and on which host.}
For t5b, $S=\num{2.43}$ measured against $F=\num{1.25}$. On the idle host of
the table above, \eqref{eq:condition} says the code loses at one, two and four
threads and wins at six. At the five contended $\Sigma_{6}$ values just
reported it says the code loses everywhere, six threads included. What is
observed is the idle-host pattern: \S\ref{sec:traffic} measures
$\num{0.53}\times$, $\num{0.62}\times$, $\num{0.78}\times$ and
$\num{1.35}\times$. Both signs therefore hold as measured; what the contended
runs withdraw is the claim that the six-thread sign is a property of the part
rather than of its occupancy.

\paragraph{These are the replay's predictions, and the four-thread one is
superseded.} Both terms of $S=\num{2.43}$ are measured against the
\code{i2\_s} \emph{reference} kernel (\S\ref{sec:microbench}), which is the
kernel the replay of \S\ref{sec:traffic} drives. It is not the kernel
\code{llama.cpp} dispatches: \S\ref{sec:endtoend} measures that path at about
$1.6\times$ the reference's time, so the ratio the condition takes in situ is
$S_{\mathrm{eff}}=\num{2.43}/\num{1.606}=\num{1.51}$, below
$\Sigma_{4}=\num{1.86}$. The four-thread cell is therefore a loss in the replay
and a win in the deployed model, and each is right about its own kernel pair.
Every prediction in this section is the replay's unless it says otherwise.

\section{Experimental setup}

\paragraph{Whose kernel the baseline is.} \code{i2\_s} is not an upstream
\code{ggml} type. \code{ggml-org/llama.cpp} carries \code{TQ1\_0} and
\code{TQ2\_0} as type numbers 34 and 35 and nothing at 36; \code{GGML\_TYPE\_I2\_S
= 36} is declared in the \code{llama.cpp} fork that \code{microsoft/BitNet}
pins as \code{3rdparty/llama.cpp}, at the commit this work builds against. Every
comparison in this paper is therefore against that fork's kernel, not against
\code{ggml}'s, and the two ternary kernels \code{ggml} does carry are examined
separately in \S\ref{sec:tq10}. The distinction matters for where a reader
looks and for where a defect would be reported.

\begin{center}
\small
\begin{tabular}{@{}l>{\raggedright\arraybackslash}p{0.62\linewidth}@{}}
\toprule
CPU & AMD Ryzen 5 3600 (Zen 2), 6 cores / 12 threads \\
\code{/proc/cpuinfo} & \code{avx2}, \code{fma}; \textbf{no AVX-512, no VNNI} \\
Caches & L1d \SI{32}{\kibi\byte}/core, L2 \SI{512}{\kibi\byte}/core, L3 \SI{32}{\mebi\byte} in two CCX slices \\
Memory & \SI{62}{\gibi\byte} \\
Compiler & gcc 13.3.0, \code{-O3 -mavx2 -mfma -march=native} \\
Model & \code{microsoft/bitnet-b1.58-2B-4T}, \num{2084044800} ternary weights, 210 tensors \\
Second model & \code{jpacifico/Aramis-2B-BitNet-b1.58-i2s-GGUF}, independent fine-tune \\
Harnesses & \code{llama-bench}, \code{llama-batched-bench}, \code{llama-perplexity} \\
\bottomrule
\end{tabular}
\end{center}

The host is shared with an unrelated production workload whose load average
rarely falls below \num{1.8}; within-arm spread ranges from \SI{2}{\percent} to
\SI{146}{\percent} with load, comparable to the effect under study. Every
comparison is paired with the arm order alternated, and headline figures are
replicated across five independent invocations of the upstream benchmark.

\section{Results}\label{sec:results}

\subsection{Correctness}\label{sec:correct}

\paragraph{Unit level.} \num{184228} assertions across five suites, zero
failures; of these the \num{177734} covering the format and both kernel variants
are in the public repository and run from a clone with a compiler alone
(\code{make test}). The other three suites cover parts of the wider project not
published here. Kernels are checked against scalar references derived independently
from the specification rather than from the implemented identities, so a shared
derivation error cannot pass. Coverage includes exhaustive verification of the
four magic-division identities over their whole legal range, deliberate 32-bit
wraparound driven \num{4.96} times past $2^{31}$, extreme activations $\pm 127$,
and every block and tail boundary.

\paragraph{Model level, argmax.} Nine prompts across two models, greedy
decoding, one binary, one seed: identical output by \code{diff}. A negative
control confirms the comparison can fail.

\paragraph{Model level, distribution.} Perplexity over 40 chunks of 512
tokens of WikiText-2, taken from \code{llama.cpp}'s CI copy:
\[
\mathrm{PPL}_{\code{i2\_s}}=\num{96.8243}\pm\num{3.55673},\qquad
\mathrm{PPL}_{\code{t5b}}=\num{96.8243}\pm\num{3.55673}.
\]
The \emph{absolute} figure is not this checkpoint's, and nothing here claims it
is: \code{llama.cpp} builds this architecture's feed-forward block with
\code{SiLU} where the checkpoint declares squared ReLU, and under the declared
activation the same forty chunks give \num{16.7482} in both arms
(\S\ref{sec:limits}). The claim this paragraph makes is the \emph{identity},
and the identity holds under both activations, because the two arms have always
run whichever graph the loader built.
Not only the final estimates but all forty running estimates agree exactly, from
$[1]\,\num{47.2744}$ to $[40]\,\num{96.8243}$; instrumentation confirms the arms
differed (\num{51660} calls through the new path against zero). This compares a
floating-point aggregate over $\num{2.05e4}$ forward passes rather than a
decoded string.

\paragraph{How much that assertion count is worth.} An assertion count measures
effort, not power, so the suite's power is measured directly by mutation:
twenty-one defects are injected one at a time into the kernel and its header ---
each of the four magic multipliers perturbed, the digit multiplication changed
from $3y$ to $5y$, the odd-byte view shifted by seven instead of eight, the
low-byte mask narrowed, the byte-wise digit subtraction made word-wise, two
digit planes given each other's activations, the packer's digit weights
perturbed, the radix changed,
the digit count reduced, the accumulator fold raised from 12 to 13 and to
the reference kernel's 32, and seven aimed at the column-blocked GEMM path specifically, which
shares no code with the per-row one and carries the \code{pp512} result.
\textbf{Eighteen are killed.}

Two of the seven are worth their own sentences, and they are the pair that
bounds the exactness argument. The proposition of
\S\ref{sec:kernel} gives a \emph{sufficient} exactness range for a magic
multiplier, and it is conservative: replacing 811 by 812 is guaranteed only
below $x = 198$ but in fact first fails at 323, above every packed byte, so it
is not a defect at all. Replacing it by 813 first fails at $x = 242$ --- exactly
the largest legal packed byte --- so it is a defect visible on one input value
in 243. The suite kills it.

The remaining three survive, and checking rather than filing them is what makes
the result meaningful: all three are \emph{exactly equivalent} on every
reachable input, proved exhaustively rather than argued. Replacing the magic multiplier
811 by 810 leaves $\lfloor x/81 \rfloor$ unchanged on all 256 byte values;
the 812 mutant just discussed is unchanged for the same reason; and replacing the
byte-wise digit subtraction by a word-wise one changes nothing on any of the
\num{65536} word lanes, because no low byte ever borrows --- the invariant
\S\ref{sec:accum} asserts, confirmed independently of the text that asserts it. The script carries both proofs and fails if either mutant is ever
\emph{killed}, since that would mean the kernel had lost the property the proof
rests on.

\paragraph{Code 3.} Absent: 0 of \num{2084044800} weights across all 210
tensors, so~\eqref{eq:shift} holds model-wide.

\paragraph{A defect in one of the two baselines.} With activations uniformly at
$\pm 127$, \code{llama.cpp}'s \code{i2\_s} \emph{vec\_dot} kernel --- the
reference of \S\ref{sec:correct}, not the \code{llamafile\_sgemm} path a real
run dispatches --- disagrees with exact arithmetic on
\textbf{\num{11998} of \num{12000}} rows at $K = 6912$, every discrepancy a
multiple of $2^{16}$ --- the overflow \eqref{eq:fold} forbids by construction in
our kernel. The prediction that overflow needs $n_b = K/128 \ge 32$, i.e.
$K \ge 4096$, is confirmed on both sides: \textbf{0 of \num{108000}} rows at
$K = 2560$ wrap under the same activations.

\textbf{The shipping path does not have this defect, and reading it is what
established that.} \code{tinyBLAS\_I2S\_AVX} folds \code{int16} into
\code{int32} once per 128-weight block, inside its block loop, where the
reference kernel accumulates 32 blocks first (\code{group32\_num = nb / 32}).
There is therefore no 32-block group to overflow on the path that runs, and the
finding above is a statement about the reference kernel alone. An earlier
version of this paragraph said ``\code{llama.cpp}'s \code{i2\_s} AVX2 kernel''
without distinguishing the two, which reads as the stronger claim.

The two exceptions are the bound speaking, not noise, and they change how the
claim must be phrased. The margin is $256 \times 127 = \num{32512}$ against
\num{32767}, or \SI{0.78}{\percent}, so whether a row overflows turns on its own
code mean and a row slightly below average clears the ceiling. At \num{12000}
rows that tail is visible, so the correct statement is statistical rather than
absolute: direction and magnitude hold exactly, ``must'' does not.

Whether ordinary activations reach that state is not established; the mechanism
is. On random \code{int8} the baseline is exact on all \num{42000} rows tried,
which is why the two kernels agree wherever it matters.

\subsection{Kernel microbenchmark}\label{sec:microbench}
One core, weights resident in L2 so transport cannot bind:

\begin{center}
\begin{tabular}{lrrrrrrr}
\toprule
kernel & GMAC/s & instr & mul & spills & weights & ceiling & of it \\
\midrule
\code{i2\_s} (the fork's) & 81.93 & 21 & 4 & 0 & 128 & 130.9 & \SI{63}{\percent} \\
t10 (word)              & 30.04 & 50.25 & 19.25 & 8.25 & 160 & 34.6 & \SI{87}{\percent} \\
\textbf{t5b (byte)}     & \textbf{33.74} & 47 & 13 & 0 & 160 & 50.5 & \SI{67}{\percent} \\
\bottomrule
\end{tabular}
\end{center}

The ceiling is $\pi\cdot(\text{weights}/\text{mul})$. The \emph{instr} column
counts every instruction in the steady-state loop body, the loop branch
included. Where \S\ref{sec:limits} and \S\ref{sec:shufdig} speak of \emph{vector}
operations they mean the vector subset of that column, which is smaller: 17 of
the baseline's 21, and 43 of t5b's 47 as built with \code{-march=native}
(\code{results/microarch\_model.txt},
\code{results/shufdig\_measurement.txt}). t5b initially measured
\num{24.36}~GMAC/s --- slower than t10 despite fewer multiplies --- which no
end-to-end benchmark could diagnose; disassembly showed 44 stack references in a
425-instruction function. One block per iteration with each plane contracted at
the moment it exists gave 47 instructions, zero spills, \num{33.74}~GMAC/s,
hence $S=\num{81.93}/\num{33.74}=\num{2.43}$.

\subsection{Weight traffic of one token}\label{sec:traffic}
Replaying one token's matrix operations over the model's real shapes and the
full weight footprint of both arms together, \SI{940}{\mega\byte}, of which the
\SI{521}{\mega\byte} of \code{i2\_s} codes is what a single arm streams, best of
five, arms rotated:

\begin{center}
\begin{tabular}{rrrrr}
\toprule
threads & \code{i2\_s} & t10 & t5b & t5b/\code{i2\_s} \\
\midrule
1 & \SI{34.57}{\milli\second} & 72.73 & 64.70 & 0.534 \\
2 & 20.29 & 36.84 & 33.01 & 0.615 \\
4 & 13.19 & 18.82 & 16.90 & 0.780 \\
6 & 15.15 & 14.49 & \textbf{11.18} & \textbf{1.354} \\
\bottomrule
\end{tabular}
\end{center}

The crossing matches~\eqref{eq:condition} with $S=\num{2.43}$ and the $\Sigma$
of \S\ref{sec:condition}. Across all eight paired six-thread observations the
median is $1.12\times$, seven of eight favouring t5b; best against best, noting
\code{i2\_s} does not improve past four threads, is \SI{13.00}{\milli\second}
against \SI{11.71}{\milli\second}, or $\num{1.110}\times$.

\paragraph{The six-thread cell is above the model's own ceiling, and the excess
is not a density effect.} \eqref{eq:condition} caps the gain at
$\min(F,\Sigma/S)$, which at six threads is
$\min(\num{1.25},\ \num{3.03}/\num{2.43})=\num{1.248}$, or \num{1.244} if
$F$ is the byte ratio this replay actually streams. The row's
$\num{1.355}\times$ stands \SI{8.6}{\percent} above the first and
\SI{8.9}{\percent} above the second, so taking the achieved ratio does not
rescue it --- it lowers the ceiling. \emph{Nothing} about weight traffic can:
$\min(F,\Sigma/S)\le F$ at every thread count, and reproducing
$\num{1.355}\times$ from bytes alone would take \num{1.476} bits per weight,
below this checkpoint's own order-zero entropy of \eqref{eq:entropy}. No
lossless re-encoding of these weights reaches it.

The two figures are not the same quantity. $\min(F,\Sigma/S)$ is a quotient of
two \emph{roofline} times, each arm at the larger of its arithmetic and
transport bounds; the row is a quotient of two measurements. At one, two and
four threads neither arm sits more than \SI{7}{\percent} above its own floor
and the modelled ratio tracks the measured one to within \SI{9}{\percent} --- at
four threads \num{0.767} against \num{0.780}. At six it does not:
\code{i2\_s}'s transport floor is \SI{13.41}{\milli\second}, essentially
unchanged from \SI{13.49}{} at four because $B$ is flat there, and the
\SI{15.15}{\milli\second} in this row stands \SI{13.0}{\percent} above it
while its t5b counterpart stands \SI{3.7}{\percent} above t5b's. That is the
widest divergence in the eight paired observations. Taking the medians instead
gives $\num{1.244}\times\num{1.064}/\num{1.207}=\num{1.097}$ against a
measured median of \num{1.120}, and best-against-best gives
$\num{1.110}\times$. \eqref{eq:condition} predicts the \emph{crossing}, and
the crossing is where it says; the size of the six-thread win is one draw's
slack above a floor, not a quantity the model claims to give.

\subsection{End to end in \code{llama.cpp}}\label{sec:endtoend}

\begin{center}
\begin{tabular}{lrrr}
\toprule
 & size & pp512 (t/s) & tg128 (t/s) \\
\midrule
\code{i2\_s} & \SI{1.10}{\gibi\byte} & $119.65\pm 2.27$ & $22.56\pm 1.26$ \\
\textbf{t5b} & \textbf{\SI{1.00}{\gibi\byte}} & $\mathbf{135.45\pm 1.21}$ & $\mathbf{25.78\pm 0.87}$ \\
ratio & 0.909 & \textbf{1.119} & \textbf{1.133} \\
\bottomrule
\end{tabular}
\end{center}

The throughput columns are one invocation's \code{llama-bench} output; the
ratio row is the \emph{median} over five, which is the figure quoted
throughout. Faster in five of five on both measures, with per-invocation ratios
$1.096$, $1.112$, $1.119$, $1.253$, $1.132$ (prompt) and $1.099$, $1.133$,
$1.123$, $1.242$, $1.143$ (generation) --- a spread wide enough that quoting any
single one of them as the result, which an earlier draft did, is a choice rather
than a measurement.

\paragraph{The four-thread result contradicted \S\ref{sec:condition}'s
prediction; the contradiction is now resolved, and the prediction was right.}
From $S=\num{2.43}$ and $\Sigma_4=\num{1.86}$, \eqref{eq:condition} predicts a
loss at four threads and the isolated replay of \S\ref{sec:traffic} delivers one
($0.780\times$); the model gains $1.119\times$, the median over five
invocations. Instrumenting the glue with a
cycle counter --- \code{rdtsc} at the invariant TSC rate of
\SI{3.599978}{\giga\hertz}, not the core clock, with \code{llama-bench}'s untimed
warmup removed by differencing $r=2$ against $r=10$ --- gives the matmul time
directly for one arm and by difference for the other, since everything else a
token does is identical between them:

\begin{center}
\begin{tabular}{lrrr}
\toprule
kernel & in situ & replay & ratio \\
\midrule
t5b & \SI{16.70}{\milli\second}/token & \SI{16.90}{\milli\second} & \textbf{0.988} \\
\code{i2\_s} & \SI{21.18}{\milli\second}/token & \SI{13.19}{\milli\second} & \textbf{1.606} \\
\bottomrule
\end{tabular}
\end{center}

The replay predicts \emph{our} kernel to \SI{1.2}{\percent} and mispredicts
\code{llama.cpp}'s \code{i2\_s} path by \SI{61}{\percent}. The effective in-situ
ratio is $S_{\mathrm{eff}} = \num{2.43}/\num{1.606} = \num{1.51}$ against
$\Sigma_4 = \num{1.86}$, so $S<\Sigma$ \emph{holds} at four threads and
\eqref{eq:condition} predicts the observed gain rather than contradicting it.

What does not survive is the use of the reference \code{i2\_s} kernel as a
stand-in for \code{llama.cpp}'s real performance: it is $1.6\times$ faster than
what the model runs, so every ratio computed against it --- the \num{2.43}
included --- flatters the baseline. No end-to-end number changes. The correction is
\emph{the baseline's}, and the table above says so: the replay predicts t5b's
matmul time to \SI{1.2}{\percent} and misses \code{i2\_s}'s by
$1.606\times$. In share terms, matmuls are \SI{49.5}{\percent} of a baseline
token where the replay implied \SI{30.8}{\percent}; for t5b the replay's
\SI{44.1}{\percent} was already right against \SI{43.6}{\percent} measured.
Stating this as a single unlabelled figure --- as an earlier draft did, taking
t5b's in-situ share against \code{i2\_s}'s replay share --- compares two arms
and understates how arm-specific the error is. Why the in-situ path costs $1.6\times$ its reference is not
established; the dispatch, the accumulator fold and the per-column
post-processing are candidates, and isolating them would mean instrumenting the
control arm.

\paragraph{The \num{1.606} is weaker than the table makes it look, and the
integration now carries the fix.} It is not measured but derived, from two
\code{llama-bench} runs taken separately on a host whose load ranged from 2 to
10 (\SI{4.905}{\second} against \SI{5.479}{\second} per repetition), so every
fluctuation \emph{between} those runs lands entirely in the \code{i2\_s} matmul
figure. The integration therefore probes \textbf{both} arms at the same
dispatch site in the same expression, so
the ratio becomes a quotient of two measured quantities and whatever the
dispatch costs is common to both and divides out.

The probe's cost was measured rather than assumed, and bounds at
\SI{0.18}{\percent} of a matmul call for the entry/exit pair, so no version of
it endangers the measurement. One fact in that null result still carries weight:
the original shared counter cost \num{4.7} times as much at six threads as at
one, because every \code{ggml} worker serialised on the same cache line, so its
error grew along the very axis this section compares on and it was replaced by a
per-thread counter.

And it has now been run~(\code{results/insitu\_symmetric.txt}). The measurement needed a patched
\code{llama.cpp}, which had meant a container this host does not grant; building
it natively instead --- \code{cmake} and \code{gcc}, no daemon, no privileges
--- removed that obstacle. \code{insitu\_measure.sh} interleaves the two
arms one repetition at a time rather than running them as two blocks, so a load
excursion reaches both at comparable rates, and forms the ratio within each
round.

Nine interleaved rounds at four threads, host load 3--5:

\begin{center}
\begin{tabular}{lrr}
\toprule
 & derived (\S\ref{sec:endtoend}, earlier) & measured \\
\midrule
\code{i2\_s} matmul, ms/token & \num{21.18} & \num{21.67} \\
t5b matmul, ms/token          & \num{16.70} & \num{16.88} \\
ratio                          & \num{1.268} & \textbf{\num{1.257}} (median) \\
rounds with \code{i2\_s} slower & ---        & 9 of 9 \\
\bottomrule
\end{tabular}
\end{center}

\textbf{The derived figures survive.} The difference method was the right thing
to object to and it gave the right answer: \SI{2.3}{\percent} on the
\code{i2\_s} time, \SI{1.1}{\percent} on t5b, \SI{0.9}{\percent} on the ratio,
with the direct measurement never once favouring the other arm across nine
rounds. What the objection bought is not a corrected number but a number that no
longer rests on the assumption that everything outside the matmuls costs the
same in both arms.

The four-thread resolution survives with them. That resolution was
built on the \num{1.606}, so it has to be recomputed on the figure that replaced
it. Against the replay's \SI{13.19}{\milli\second} the measured
\SI{21.67}{\milli\second} is a factor $\num{1.643}$, giving
$S_{\mathrm{eff}} = \num{2.43}/\num{1.643} = \num{1.48}$ against
$\Sigma_{4}=\num{1.86}$. The inequality $S<\Sigma$ holds at four threads on the
measured factor as it did on the derived one, and by a wider margin.

\paragraph{Where that time goes.} Two further probes inside
\code{tinyBLAS\_I2S\_AVX}, one around the contraction loop and one around the
per-column post-processing, separate what the dispatch-site probe cannot.
Dispatch is the remainder, measured by difference, because a probe around the
call would be the call. Five runs:

\begin{center}
\begin{tabular}{lr}
\toprule
part of the \code{i2\_s} matmul & share \\
\midrule
contraction loop        & \SI{94.4}{\percent} \\
post-processing         & \SI{3.1}{\percent} \\
dispatch (by difference) & \SI{2.5}{\percent} \\
\bottomrule
\end{tabular}
\end{center}

All three candidates are now answered. Dispatch and post-processing
together are \SI{5.6}{\percent}, so neither carries the difference: the
\num{1.257} is a statement about the inner loop's own arithmetic and memory
traffic. The third, the accumulator fold, needs no probe at all: as
\S\ref{sec:correct} establishes from the source, \code{tinyBLAS\_I2S\_AVX} folds
once per 128-weight block and has no 32-block group. The fold is therefore not a
periodic cost that occasionally interrupts the contraction; it is part of every
block, and it is inside the \SI{94.4}{\percent}.

That also sharpens why the overflow of
\S\ref{sec:correct} happens where it does. The kernel feeds
\code{vpmaddubsw} with raw 2-bit codes against \code{int8} activations, so an
\code{int16} lane takes at most $2 \cdot 3 \cdot 127 = 762$ per instruction and
$4 \cdot 762 = 3048$ per block: \textbf{ten blocks may be accumulated, not
thirty-two}. The reference kernel's 32 is three times its own bound, which is
the overflow; the dispatched kernel's 1 is a tenth of it.

The nine omitted folds cost \SI{7.5}{\percent} of the dispatched path's
matmul time. A tree identical but for accumulating ten blocks before folding
measures $9.49 \times 10^{9}$ cycles against $10.25 \times 10^{9}$, three runs
each, with \emph{identical} perplexity --- \num{108.6810} on four chunks of
WikiText-2 in both. That is one measured component of the \SI{25}{\percent} by
which the dispatched path trails t5b, not an explanation of it, and it is not a
proposal: the \code{i2\_s} path is the control arm of every comparison here, the
modified kernel lives in its own tree, and no reported ratio is built from it.

These probes are compiled out by default, and that is not tidiness. They fire
once per $\mathit{RM} \times \mathit{RN}$ tile rather than once per call, and
they instrument one arm only, so leaving them in inflates the \code{i2\_s} side
of the ratio --- measured at \num{1.379} against \num{1.257} on the same host, a
\SI{10}{\percent} artefact of the instrument. The decomposition and the ratio
come from two builds by design.

One thing this still does not settle: it is one host, the same shared machine as
everything else in \S\ref{sec:results}.

\paragraph{Prompt length.} $1.106$, $1.152$, $1.180$, $1.113$ at 64, 256, 1024
and 2048 tokens: the advantage does not decay with length. We had predicted
decay on the grounds that a prompt pass is arithmetic-bound; at these shapes the
weights still stream from memory, so the byte saving remains live and the
prediction was wrong in its premise.

\paragraph{Concurrency.} Total throughput at 1, 2, 4, 8 parallel sequences:
$1.041$, $1.041$, $1.027$, $1.097$. One invocation per model at each width, no
repetition, and the $N=4$ prompt cell is a recorded outlier
(\code{results/inference\_t5b.txt}), so this is a shape and not a trend: the
series is flat, then falls, then rises, and only the last point would support
the reading that the advantage \emph{grows}. The mechanism that would make it
grow is real --- $N$ concurrent sequences turn generation into a
matrix--matrix product with $N$ columns, across which the decode amortises ---
and \S\ref{sec:shufdig} measures that amortisation directly and finds it
complete by $N=8$. These four points do not establish it.

\paragraph{Across checkpoints.} On the independent fine-tune: identical output
on three prompts, $1.120\times$ prompt and $1.075\times$ generation.

\paragraph{Across depth.} Repeating the block stack yields a
\num{4.50}\,B-parameter, 60-layer graph \code{llama.cpp} loads and runs with
real attention and KV cache: $1.165\times$ and $1.129\times$, both margins clear
of their error bars. That artefact is \emph{not a model} --- the same thirty
blocks run twice --- and no quality claim is made from it.

\subsection{Two ratios a reader will ask about immediately}\label{sec:ratios}

\paragraph{Why \SI{19.6}{\percent} fewer weight bytes give \SI{13.3}{\percent}
more throughput.} The \SI{13.3}{\percent} is the median generation ratio of
\S\ref{sec:endtoend} over five invocations, which is the figure quoted
throughout. Both halves of the question below come from one of those runs, whose
own end-to-end ratio is \num{1.117}. The
\code{rdtsc} counter of \S\ref{sec:endtoend} times the matmuls of a
\code{llama-bench} \code{tg128} invocation at four threads, and the same
invocation's wall clock gives the throughput, so the decomposition and the ratio
it explains are not assembled from different runs or different arms. Per token:

\begin{center}
\begin{tabular}{lrrrr}
\toprule
 & wall & matmuls & matmul share & the rest \\
\midrule
\code{i2\_s}  & \SI{42.80}{\milli\second} & \SI{21.18}{\milli\second} & \SI{49.5}{\percent} & \SI{21.6}{\milli\second} \\
\textbf{t5b} & \textbf{\SI{38.32}{\milli\second}} & \textbf{\SI{16.70}{\milli\second}} & \SI{43.6}{\percent} & \SI{21.6}{\milli\second} \\
\bottomrule
\end{tabular}
\end{center}

\textbf{That table is a budget, not evidence, and an earlier draft of this
paragraph presented it as evidence.} Only the t5b arm carries the counter;
\code{i2\_s}'s matmul time is \emph{defined} as its wall clock minus the
non-matmul time measured in the other arm. The last column is therefore
identical between the arms by construction, and the ``prediction'' that
\SI{49.5}{\percent} of a baseline token, reduced by \SI{21.1}{\percent},
gives $\num{1.117}$ end to end is the identity
$T_{A}/\bigl(T_{A}-(M_{A}-M_{B})\bigr)=T_{A}/T_{B}$ --- it reproduces the
measured ratio exactly, for any inputs, and cannot disagree with it. Quoting
that agreement as a cross-check was a check that could not fail.

The matmul ratio \emph{is} measured, elsewhere and on both arms: the symmetric
probe of \S\ref{sec:endtoend} times \code{i2\_s} and t5b at the same dispatch
site over nine interleaved rounds and gives a median of $\num{1.257}$, with
\code{i2\_s} slower in 9 of 9 (\code{results/insitu\_symmetric.txt}). The
budget above is worth keeping for what it does show --- that the untouched
\SI{21.6}{\milli\second} is the larger half of a token, so a format acting on
the smaller half cannot do better than it does --- and not for a corroboration
it cannot supply.

Two figures in that chain deserve their arm stated, because taking the wrong one
is the mistake this paragraph previously made. \SI{43.6}{\percent} is the matmul
share \emph{after} the saving, so it answers ``what does t5b spend its time on''
and not ``why is it faster''; the second question takes \SI{49.5}{\percent}. And
an independent instrument agrees: the byte ratio
$\num{521011200}/\num{418775040}=\num{1.244}$ against the clock's \num{1.268},
two measurements of the same quantity that share no apparatus, differing by
\SI{1.9}{\percent}.

\paragraph{What is in the untouched half, and who does act on it.} The
\SI{21.6}{\milli\second} is identical in both arms by construction: this work
changes how ternary \emph{weights} are stored and decoded, and nothing else in
the graph. Attention scores, the norms, the sampler and --- the largest
memory-resident item among them --- the key/value cache are all outside it. That
is the honest reason a weight format and a cache format compose rather than
compete, and it is worth stating precisely, because the tempting shorter
argument is wrong: t5b does \emph{not} leave the attention tensors alone. All
\num{210} ternary tensors are re-encoded, \code{attn\_q} among them
(\code{results/tensor\_structure.txt}). What it leaves alone is the cache, which
holds activations and never weights.

Hariri et al.~\cite{spectralkv} act on exactly that half. They show the key
projections carry larger spectral and Frobenius norms than the value
projections, and allocate bits accordingly --- four to keys, two to values ---
retaining up to \SI{98.3}{\percent} of the accuracy of a uniform allocation at
lower memory. The two results are orthogonal in the strict sense that they share
no tensor: a reader combining them would shrink the weight side by the
\SI{102.2}{\mega\byte} of \S\ref{sec:results} and the cache side by their
allocation, and neither measurement here predicts the other. We have not run the
combination, and say so rather than implying a composition we did not test.

The standalone replay of \S\ref{sec:traffic} puts the baseline's matmuls at
\SI{13.19}{\milli\second} against the \SI{21.18}{\milli\second} measured in
situ, for the reason \S\ref{sec:endtoend} gives: it drives
\code{bitnet\_\allowbreak vec\_\allowbreak dot\_\allowbreak i2\_\allowbreak i8\_\allowbreak s\_\allowbreak reference}
and not the path \code{llama.cpp} dispatches, so it is a lower bound on what the
matmuls cost. Read as bandwidth, the weight stream moves
\SI{521.0}{\mega\byte} in the \SI{21.18}{\milli\second} the baseline spends in
the matmuls --- \SI{24.6}{\giga\byte\per\second}, against \SI{25.1}{} for t5b
over its \SI{418.8}{\mega\byte} and \SI{16.70}{\milli\second}. Equal rates in
both arms, at \SI{64}{\percent} of the \SI{38.62}{\giga\byte\per\second} this
machine delivers at four threads: the format changes how many bytes cross, not
how fast they cross, and neither arm is at the wall.

\paragraph{Why the file shrinks by only \SI{9}{\percent}.} The ternary tensors
are \SI{521.0}{\mega\byte} of a \SI{1187.8}{\mega\byte} file. One f16 tensor ---
the 128k-vocabulary embedding --- is \SI{656.7}{\mega\byte}, \SI{55.3}{\percent}
of the file, and does not change. The \SI{102.2}{\mega\byte} saved is
\SI{19.6}{\percent} of what the format touches and \SI{8.6}{\percent} of the
file. The ternary share, and therefore the file saving, rises with model depth
at fixed vocabulary.

\subsection{A negative result on row tiling}
Tiling four weight rows against one activation vector measures $1.19\times$ in
cache and $0.52\times$ from memory. A tile of $R$ rows by $C$ columns loads each
weight once and spends it on $RC$ contractions; a matrix--vector product has
$C=1$, so tiling rows buys no reuse and merely reorders loads. Upstream draws
the same line structurally: its \code{ggml\_gemv\_i2\_i8\_s} does not tile,
while the $4\times 4$ tile lives in \code{ggml\_gemm\_i2\_i8\_s}.

\subsection{A cheaper decode, and the width at which it stops
mattering}\label{sec:shufdig}
Because $q_3=\lfloor x/27\rfloor\le 9$ for every byte, it is a nibble, and
\code{vpshufb} can index it. The idiom --- once a quantity is provably below
sixteen, stop computing it and shuffle it --- is standard and is not claimed
here~\cite{mulapshufb}; what is measured below is a smaller decode, not a new
technique. It differs from the table-based ternary kernels named in \S1 in the way that matters for their published cost: their
shuffle is rebuilt per activation block and \emph{replaces} the multiply--add,
whereas these three tables are loop-invariant constants and the
\code{vpmaddubsw} contraction is untouched. Three constant 16-entry tables then yield
$d_4=\lfloor q_3/3\rfloor$, $d_3=q_3\bmod 3$ and $3q_3$ from $q_3$ alone, so the
multiplier $M_4=\num{811}$ and both of its \code{vpmulhuw} disappear. The packed
representation is untouched --- same five digits, same \num{1.60755}
bits/weight, same bytes --- so the variant is a build option
(\code{-DTERNARY\_T5B\_SHUFDIG}) and not a second format. It is \emph{off} by
default, and every figure elsewhere in this paper --- every $S$, every rate,
every throughput --- is from the build without it. It is bit-identical: the t5b
kernel suite --- \num{97529} assertions, the smallest of the three counts this
paper quotes and the one the mutation harness of \S\ref{sec:correct} uses as its
control, not the \num{184228} across five suites --- passes against it, and the
two decoders
agree on all 256 byte values including the thirteen foreign ones.

Under the build's own flags the loop falls from 43 to 37 vector operations and
from 13 to 11 multiply-port operations.\footnote{The count depends on
\code{-march=native}, which the \code{Makefile} sets: without it the same source
gives 42 and 36. The \num{0.2625} vector operations per weight quoted in
\S\ref{sec:limits} is $42/160$, the count \emph{without} that flag, which is
also what \code{microarch\_model.sh} compiles.} Fifteen interleaved rounds, with the
untouched \code{i2\_s} kernel measured in the same binaries as a control, give a
median of $\num{1.2254}\times$ on the contraction, faster in 14 of 15 rounds
--- the exception reading $\num{0.6086}$, an excursion eight times the control's
spread, caused by a second benchmark started on the same pinned core while the
run was going and left in because excluding it would raise the median to
$\num{1.2257}$ and the significance to $p=\num{0.000122}$ ---
\footnote{\code{shufdig\_ab.sh} reproduces every figure in this subsection. It
compiles both arms from the one file, asserts that the arm without the
definition is byte-identical to the shipped kernel, chooses its core by sampling
\code{/proc/stat} rather than by convention --- a busy SMT sibling costs as much
as a busy core --- and prints the control beside the effect, so a reader can see
the noise floor instead of being told it.}
$p=\num{0.00098}$ by an exact two-sided sign test; the control's median is
$\num{1.0037}\times$ with a round-to-round range of \SI{5}{\percent}, which is
this measurement's noise floor. $S$ falls from \num{2.43} to \num{1.98}.

\emph{No cell of the surplus table changes.} The $S$ that moves here is the
replay's: both its terms are measured against the \code{i2\_s} \emph{reference}
kernel of \S\ref{sec:microbench}, not against the path \code{llama.cpp}
dispatches, so what follows is a statement about the same frame as
\S\ref{sec:condition}'s table and not about the in-situ ratio of
\S\ref{sec:endtoend}. In that frame \eqref{eq:condition} failed at one
to four threads and still fails; it held at six and still holds. What changes is
the margin at the top of that range: the four-thread cell moves from
\SI{31}{\percent} short to \SI{6.5}{\percent} short. The improvement lands
exactly where the format was losing and does nothing where it was already
winning, because at six threads the packed kernel already sits at the memory
knee and additional arithmetic cannot be spent.

It also does nothing for prompt processing, and the reason is structural rather
than incidental. \code{ternary\_t5b\_gemm\_avx2} decodes each block once and
spreads it over $C$ activation columns, so the decode is amortised over the
strip width. Measured across widths at $K=\num{6912}$, the gain is
$\num{1.1608}\times$ at $C=1$ and $\num{1.1363}\times$ at $C=4$, both in 9 of 9
rounds, and by $C=8$ both arms saturate at the same \num{78}~GMAC/s
with 3 to 5 wins in 9 --- a coin flip. Past that width the cost is the
contraction and the activation stream, and the decode, which is all this change
touches, has been amortised away. Only \code{tg128} can move.

Finally, the measurement exceeds what \S\ref{sec:limits}'s own model predicts:
$43/37=\num{1.1622}$ on total vector operations, $13/11=\num{1.1818}$ on the
multiply port, against \num{1.2254} measured. The instruction mix shows the
multiply port falling from 13 to 11 while the shift-and-shuffle pair stays at
exactly 9 --- three shifts traded for three shuffles --- so the distribution
across ports flattens. On the \num{18.54} cycles per block this host's measured
rate and clock imply (\S\ref{sec:microarch}) the loop sustains \num{2.45} vector
operations per cycle against the baseline's \num{2.32}. The kernel header and
\code{results/shufdig\_measurement.txt} take \num{19.7} cycles for the same
block instead, so the latter prints \num{2.18} rising to \num{2.30}: the two
cycle counts differ by \SI{6}{\percent} and the ratios they give agree to
\SI{0.1}{\percent}.
All four figures are above the two per cycle that a part cracking
every 256-bit integer operation into two 128-bit micro-operations should
sustain, which is itself a reason to treat these absolute figures as indicative
and the ratio as the measurement. Both loops spill zero times, so register pressure
is not the mechanism. The flatter distribution is \emph{consistent} with the
residual \SI{5}{\percent} and is not isolated by this experiment; what the
experiment does establish is that removing operations is not the only lever, and
that the model of \S\ref{sec:limits} is therefore incomplete.

\subsection{Against \code{TQ1\_0}, the nearest prior art}\label{sec:tq10}

\S\ref{sec:priorart} records that the base-3 packing is \code{TQ1\_0}'s, and
that Spectra~1.1~\cite{spectra} publishes the same $p=8$, $k=5$ construction
with a correctness theorem. This subsection measures against \code{TQ1\_0},
which is the form of it that ships in the codebase integrated into; no run
against Spectra's own kernels was made. The measurement is in \code{results/tq1\_0\_comparison.txt}. The comparison file is produced by re-encoding
the same ternary values into \code{TQ1\_0}'s block layout, with the per-tensor
\code{i2\_s} scale written into every block's f16 field, so the two files carry
identical weights.

\begin{center}
\begin{tabular}{lrrr}
\toprule
 & file & ternary payload & achieved bits/weight \\
\midrule
\code{i2\_s}      & \num{1187801280} & \num{521011200} & \num{2.00000} \\
\code{TQ1\_0}     & \num{1106386560} & \num{439603200} & \num{1.68750} \\
\textbf{t5b}      & \textbf{\num{1085565120}} & \textbf{\num{418775040}} & \textbf{\num{1.60755}} \\
\bottomrule
\end{tabular}
\end{center}

\code{TQ1\_0}'s \num{1.6875} is 54 bytes per 256 weights: 48 for the
five-per-byte codes, 4 for a 16-element remainder at two bits, and 2 for the f16
block scale. t5b's \num{1.60755} is 32 bytes per 160 weights plus the tail waste
at $K = 6912$ and one 32-byte record per tensor.

\paragraph{On size the honest margin is small.} \SI{4.7}{\percent} of the
ternary payload and \SI{1.9}{\percent} of the file. Four fifths of the saving
this work reports against \code{i2\_s} was already present in the vendored tree
under an upstream type number, and four fifths of the remaining fifth was
present in the same pull request as \code{Q1\_3} before \code{TQ1\_0} replaced
it. \S\ref{sec:priorart} works that residual out in bytes: what is attributable
here is a block-length choice and nothing else.

\paragraph{On speed the margin is large and mostly not about the packing.}
Three rotated \code{llama-bench} invocations, all three files in each:

\begin{center}
\small
\begin{tabular}{l>{\raggedright\arraybackslash}p{0.44\linewidth}l}
\toprule
 & \code{pp512} (t/s) & \code{tg128} (t/s) \\
\midrule
\code{i2\_s}  & $122.25 \pm 1.52$, $77.06 \pm 4.49$, $118.12 \pm 3.26$ & 23.61, 15.27, 23.56 \\
\code{TQ1\_0} & $48.27 \pm 0.73$, $43.95 \pm 5.26$, $49.01 \pm 0.29$   & 21.68, 16.07, 22.27 \\
\textbf{t5b}  & $130.65 \pm 4.35$, $137.38 \pm 0.52$, $128.17 \pm 8.07$ & 22.95, 26.43, 26.43 \\
\bottomrule
\end{tabular}
\end{center}

\code{TQ1\_0} loses prompt processing by a factor of about three, and the reason
is structural rather than arithmetic: \code{i2\_s} has a
\code{llamafile\_sgemm} case upstream, t5b has one because we wrote it, and
\code{TQ1\_0} has none --- so it processes a prompt one activation column at a
time. \textbf{A \code{TQ1\_0} sgemm case, which nobody has written, would be
expected to erase that margin.} At \code{tg128}, where no format has a batched
path to exploit, all three land within about \SI{10}{\percent} and the ordering
does not survive this host's noise.

\paragraph{The depth-one form, transplanted into upstream's own kernel.}
Upstream names the depth-one AVX2 decode as an open question in its own source
(\S\ref{sec:priorart}). Transplanting the form into \code{TQ1\_0}'s kernel and
measuring it there answers that question, and the answer is no: the transplant
measures $\num{0.983}\times$ --- faster in only 7
of 54 interleaved rounds across six invocations --- at $131\to137$ vector
operations and $0\to4$ register spills. The reason does not transfer to
the format described here, and the difference is the encoding.
\code{TQ1\_0} stores $\lceil 256v/243\rceil$, a fixed-point fraction, so
its digits come out by multiplying by $3^{j}$ in \emph{byte} lanes and
taking the top bits --- no multiply-port instruction at all, and no reason
to leave byte lanes. t5b stores the integer $v\le\num{242}$, so its digits
are floor divisions, AVX2 has no 8-bit multiply, and the even/odd split
and fold that costs \code{TQ1\_0} those six operations is not an
alternative here but the only route. The instantiation is therefore real
and is not worth what a reader might assume: it is the right choice for
this encoding and the wrong one for upstream's.

\paragraph{What survives.} Against \code{TQ1\_0} on this machine, t5b is worth
\SI{4.7}{\percent} of the ternary bytes, plus the fact of being wired into the
batched path. The claim that a sub-two-bit ternary packing is itself novel does
not survive. It is withdrawn in \S\ref{sec:priorart} along with the scale
model, the depth-one extraction arithmetic, the fused digit-plane contraction
and the condition itself, and with the interleave, which \S\ref{sec:format}
concedes against both upstream formats: six constructions, the same six the
abstract names. What survives is narrower, and is the list of
\S\ref{sec:contrib}.

\section{Dead feed-forward neurons}\label{sec:dead}

\SI{2.9036}{\percent} of all weight-matrix rows are identically zero --- an
output neuron that cannot fire for any input.

\begin{center}
\begin{tabular}{lrr}
\toprule
tensor & dead rows & share \\
\midrule
$W_{\text{gate}}$ & \num{9870} / \num{207360} & \SI{4.76}{\percent} \\
$W_{\text{up}}$   & \num{9870} / \num{207360} & \SI{4.76}{\percent} \\
$W_{v}$           & 103 / \num{19200} & \SI{0.54}{\percent} \\
$W_{q}$           & 4 / \num{76800} & \SI{0.01}{\percent} \\
$W_{\text{down}},W_{k},W_{o}$ & 0 & \SI{0.00}{\percent} \\
\bottomrule
\end{tabular}
\end{center}

They concentrate at the front: \SI{43.59}{\percent} of layer 1's feed-forward
neurons, \SI{27.58}{\percent} of layer 2's, \SI{14.77}{\percent} of layer 3's,
near zero by layer 8. \textbf{In all thirty layers the dead index set of
$W_{\text{gate}}$ equals that of $W_{\text{up}}$} --- not the same cardinality,
the same indices. In a gated block, $h = W_{\text{down}}\!\left(f(W_{\text{gate}}x)\odot W_{\text{up}}x\right)$
--- this checkpoint's $f$ is a squared \textsc{relu} and not the \textsc{swiglu}
an earlier draft named here, as \S\ref{sec:limits} establishes ---
the two are combined element-wise, so a zero
gate row renders the matching up row irrelevant and conversely; training removed
them in pairs. $W_{\text{down}}$ has none, the same fact along the other axis:
there a dead neuron is a dead \emph{column}.

They survive fine-tuning exactly: against the independent fine-tune, 210 of 210
tensors have identical dead index sets, all \num{19847} rows,
\SI{100.00}{\percent}. Whatever removed them occurred in pre-training.

Skipping the rows in the kernel saves \SI{2.4380}{\percent} of weight bytes
\emph{and} the same share of multiply--accumulates, and is unconditionally
exact. Removing the neurons outright was also attempted, and the outcome is a
negative result worth stating because three prior claims about it --- two of
them ours --- were wrong.

It is expressible: \code{llama.cpp} reads \code{LLM\_KV\_FEED\_FORWARD\_LENGTH}
through \code{get\_key\_or\_arr} into a per-layer array
(\code{llama-model.cpp:1117}), so the container and the hyper-parameters carry
per-layer widths. It is \emph{not} free of loader changes, which an earlier
draft of this section asserted without checking:
\code{src/models/bitnet.cpp} allocates its four feed-forward tensors with the
global \code{n\_ff}, so a file whose layers differ in width fails at load. Four
lines change that, and with a scalar \code{feed\_forward\_length} the broadcast
makes them a no-op for every existing checkpoint.

With those four lines the pruned model loads --- \num{2.35}\,B parameters
against \num{2.41}\,B, \num{1.08}\,GiB against \num{1.10}, \num{60948480}
ternary weights and \num{15268736} bytes removed. Every surviving weight is
bit-identical and every attention tensor byte-identical; the feed-forward
block's int8 input is unchanged value for value across all \num{1642496} of
them. But it is \textbf{not} exact at the token level, and the reason is the
sub-layer RMS norm BitNet places between the element-wise product and
$W_{\text{down}}$. That norm divides by the root mean over all \num{6912}
entries, so deleting $d$ of them rescales the output by
\[
    \sqrt{\frac{K-d}{K}}\,,
\]
which in layer 1, where $d/K = \SI{43.59}{\percent}$, is \num{0.751} and not a
rounding matter. Pinning the divisor to the original width restores it to
\num{2.5} ulp of float32 --- not to zero. On three prompts the output is
identical on one and diverges on two, late, inside repetitive passages where
successive candidates are near-ties and the last bits decide. And it buys
nothing measurable: \code{pp512} $\num{119.24} \pm \num{3.11} \to
\num{109.90} \pm \num{9.48}$, \code{tg128} $\num{23.49} \pm \num{0.51} \to
\num{23.53} \pm \num{0.72}$.

The honest summary is that \SI{2.9}{\percent} of the ternary payload can be
removed, that doing so costs four lines and changes the output, and that it
makes the model no faster. We do not ship it.

\section{Limitations}\label{sec:limits}

\paragraph{The port analysis is Zen 2 specific --- but $S$ is not.}
\S\ref{sec:condition} measures one multiply-class port at \num{1.05}
instructions per cycle against \numrange{3.0}{4.0} for the cheap classes. A
reviewer ran the same benchmark on an Intel part with AVX-512 and obtained
\num{1.60} against \num{3.09}: two 256-bit multiplier ports rather than one. The
ratio \eqref{eq:condition} turns on is therefore a property of the
microarchitecture at a finer grain than the instruction set, and the surplus
table would have to be re-measured on any other part before \eqref{eq:condition}
could be applied to it.

The worry this raises is sharper than the port count alone, and it deserves its
strongest form. Zen 2 cracks every 256-bit integer vector operation into two
128-bit micro-operations issued over four floating-point pipes, retiring roughly
two vector operations per cycle; Zen 3 and later, and Ice Lake and later,
execute 256 bits natively. The packed kernel spends $42/160=\num{0.2625}$ vector
operations per weight against the baseline's $17/128=\num{0.1328}$: it is the
kernel that issues \emph{more} of them. A machine doubling the vector-op ceiling
therefore relieves the packed code preferentially, $S$ could collapse, and the
whole motivation for a roofline condition would be an artefact of one narrow
part. It does not collapse: \S\ref{sec:microarch} puts $S$ between \num{2.10}
and \num{2.45} on Zen 3/4/5, Ice Lake and Sapphire Rapids against \num{2.50} on
Zen 2, and \S\ref{sec:second} measures \num{2.11} on an Intel part.

\paragraph{One microarchitecture, and the decisive property is absent from it.}
All results are from a single Zen 2 part reporting \code{avx2} and \code{fma}
and nothing else. \code{VPDPBUSD} collapses the contraction into one instruction
while leaving the decode unchanged; by~\eqref{eq:surplus} this raises $S$ and
lowers $\Sigma$, moving~\eqref{eq:condition} against a packed code. The
\emph{Litespark} system, targeting VNNI and ARM \code{SDOT}, stores ternary
weights at eight bits for precisely this reason. \S\ref{sec:microarch}
\emph{bounds} that move at \numrange{0}{14}\,\% from a static model;
\S\ref{sec:vnnimeas} \emph{measures} it at \SI{8.8}{\percent} on hardware we do
not own. \textbf{The central result is still ours only for AVX2 without VNNI}
--- every end-to-end number in \S\ref{sec:results} comes from this one host, and
no VNNI machine has run the model. What \S\ref{sec:vnnimeas} settles is the
kernel-level ratio, not the deployed one.

\subsection{How far a model can substitute for the hardware}\label{sec:microarch}

The two objections above differ in one respect that matters: the first concerns
code that exists and can be analysed, the second concerns code that does not.

We run \code{llvm-mca}, LLVM's static pipeline simulator, over the inner loops
\code{gcc -O3 -mavx2 -mfma} actually emits for both kernels, against the
vendor-derived scheduler models for six targets. Re-deriving the loop shapes
reproduces \S\ref{sec:microbench}'s counts exactly for the baseline (21
instructions, 4 multiply-class, 128 weights) and to within one instruction for
the packed kernel (46 against 47; the 13 multiplies and 160 weights are exact),
the difference being a compiler version.

The model is calibrated before it is used. Its Zen 2 prediction is
$S=\num{2.501}$ against the $S=\num{81.93}/\num{33.74}=\num{2.428}$ measured on
this host: an error of \SI{3.0}{\percent}. The script exits non-zero should that
drift past \SI{8}{\percent}. Absolute cycle counts are optimistic, which is why only
the ratio is carried forward --- but by how much this paper cannot say
precisely, and an earlier draft overstated its own confidence here. It quoted
\num{18.01} predicted against \num{19.7} measured per 160-weight block, a
\SIrange{6}{11}{\percent} band. That \num{19.7} \emph{is} recorded --- in the
kernel's own header, which ships publicly as \code{src/ternary\_t5b.c}, and in
\code{results/shufdig\_measurement.txt} --- but it is nowhere \emph{derived},
and it disagrees with the figure this paper's own two constants imply:
\num{33.74}\,GMAC/s at
\SI{3.91}{\giga\hertz} gives $160\times\num{3.91}/\num{33.74}=\num{18.54}$
cycles, against which the model is optimistic by \SI{2.9}{\percent}. The
\num{19.7} would require a clock of \SI{4.15}{\giga\hertz}. We report the
derived figure and the discrepancy rather than choosing between them, and
\S\ref{sec:shufdig} states which of the two each vector-operations-per-cycle
figure is backed out of.

\begin{center}
\begin{tabular}{lrrr}
\toprule
target & packed, cyc/block & baseline, cyc/block & $S$ \\
\midrule
\code{znver2} (this host) & 18.01 & 5.76 & 2.501 \\
\code{znver3}             & 11.51 & 3.76 & 2.449 \\
\code{znver4}             & 11.51 & 3.76 & 2.449 \\
\code{znver5}             & 11.51 & 3.76 & 2.449 \\
\code{icelake-server}     & 14.02 & 5.34 & 2.098 \\
\code{sapphirerapids}     & 14.52 & 5.01 & 2.317 \\
\bottomrule
\end{tabular}
\end{center}

The third significant digit of that table is a property of the translation unit
and not of the target --- the public copy of
\code{results/microarch\_model.txt} prints \num{3.77} and $S=\num{2.441}$ for
\code{znver3}, \code{znver4} and \code{znver5} where this one prints \num{3.76}
and \num{2.449}, because the two \code{i2\_s} inner loops are twenty
instructions each with an identical opcode multiset and identical memory offsets
and differ only in register allocation and order, which \code{llvm-mca} prices
at \num{0.01} cycles on those three targets and at zero on the other three ---
so it costs \SI{0.33}{\percent} on $S$, and the claim the table is cited for
survives it (\code{results/microarch\_translation\_unit.txt}, which isolates
the effect to those twenty instructions and records the three zeros as its
control).

\textbf{These are modelled numbers, and the licence for printing them is that
the model was tested twice.} Once in-sample, against the host it was calibrated
on (\SI{3.0}{\percent}); once \emph{out of sample}, against a machine it had
never seen --- the reviewer run of \S\ref{sec:second} measured $S = \num{2.11}$
on an Intel part, where this table predicts \num{2.098} and \num{2.317} for the
two Intel targets. A model that reproduces its calibration set proves nothing;
one that predicts a case withheld from it has earned a table. Every figure here
is nonetheless labelled modelled, and the sections it feeds say so at each use.

Two readings would overstate this table. \code{znver3}, \code{znver4} and
\code{znver5} return \emph{identical} counts, so they are one prediction and not
three: it holds four distinct predictions, not six. And every newer target sits
at or below the Zen 2 value rather than scattered around it --- the ratio
improves by up to \SI{16}{\percent} and never degrades --- so what is supported
is one-sided, that $S$ does not run away, and not that $S$ is constant.

For VNNI the same apparatus is far weaker, and we set out why rather than
reporting the number alone. No compiler here targets VNNI usefully, so the two
VNNI loops are written by hand: in both kernels the contraction ends in a
\code{vpmaddubsw} followed by an accumulating \code{vpaddw}, and \code{VPDPBUSD}
performs exactly that pair fused, with the same unsigned-times-signed operand
convention. The substitution is argued sound --- the reduction width changes
from \code{int16} pairs to \code{int32} quadruples, leaving the horizontal sum
invariant, and it makes the accumulator bound of \S\ref{sec:accum} unnecessary
--- but that is an argument, not a test. \textbf{These loops have never been
executed.} Three caveats bound what the numbers mean.

\begin{enumerate}
\item \code{VPDPBUSD} accumulates in place with latency \numrange{10}{13}, so a
loop carrying one accumulator per contraction is bound by that latency until
unrolled. On \code{sapphirerapids} this makes the hand-written VNNI
\emph{baseline} slower than the AVX2 one it replaces --- \num{7.01} cycles
against \num{5.01} --- which is a property of the accumulator count, not of
VNNI. The comparison is therefore read at the resource bound, where an
adequately unrolled implementation lands.
\item That metric fails the calibration the dependency-bound metric passed: on
Zen 2 it gives $S=\num{1.96}$ against the measured \num{2.428}, an error of
\SI{-19}{\percent}, because the real loop is dependency-bound. Both are
reported.
\item {\sloppy The negative control did not fire. \code{llvm-mca} 23.1.0 assembles and
times \code{VPDPBUSD} for \code{-mcpu=znver2}, a part without the instruction,
and continues to under \code{-mattr=-avx512vnni,\discretionary{}{}{}-avxvnni},
assigning it latency
11 and throughput \num{1.00}. The tool cannot distinguish a legal VNNI sequence
from an illegal one and so offers no check on this subsection beyond timing. The
control is kept in place and reported failing.\par}
\end{enumerate}

Subject to all three, VNNI moves $S$ in the predicted direction --- against the
packed code, the baseline spending 4 of 17 operations in the contraction
\code{VPDPBUSD} collapses where the packed kernel spends 5 of 42 --- and moves
it by \numrange{0}{14}\,\%: $S$ reaches \num{2.08} on \code{znver4} and
\num{2.51} on \code{icelake-server} and \code{sapphirerapids}. Modest, not
decisive.

None of this is a measurement. It also cannot become one by
recompiling: no compiler contracts \code{vpmaddubsw} followed by an accumulating
\code{vpaddw} into \code{VPDPBUSD} as an idiom, so an AVX2 kernel built with
\code{-march=native} on a VNNI part contains \textbf{zero \code{vpdpbusd}} and
measures AVX2 on that part. The instruction has to be written.

\code{bench\_vnni.c} now writes the instruction: \code{i2\_s} and t5b in both
forms, the VNNI pair contracting into \code{int32} accumulators, which for t5b
also removes the fold of \S\ref{sec:accum} entirely. It verifies with
\code{objdump} at run time that its own binary contains the instruction, checks
every row against exact \code{int64} arithmetic before printing a rate, and
exits 3 rather than run on a CPU without the feature. A second build models
\code{VPDPBUSD} in AVX2 so the algorithm can be checked where the instruction
cannot execute; it passes on all 100 rows for all four kernels and deliberately
prints no timings.

\subsection{VNNI, measured}\label{sec:vnnimeas}

A reviewer ran that benchmark on an Intel Xeon reporting \code{avx512\_vnni}:
eight invocations, the ISA check passing on each, all four kernels exact against
\code{int64} arithmetic on every row. \textbf{This is the paper's largest
limitation moving from modelled to measured}, and it moves only part of the way,
so both parts are stated.

\begin{center}
\begin{tabular}{lrrr}
\toprule
 & AVX2 & VNNI & \\
\midrule
median $S$ (\code{i2\_s}/t5b) & \num{2.04} & \textbf{\num{2.22}} & $+\SI{8.8}{\percent}$ \\
\code{i2\_s} gain from VNNI & --- & --- & $1.24\times$ \\
t5b gain from VNNI          & --- & --- & $1.15\times$ \\
runs in which $S$ rose      & --- & 6 of 8 & \\
\bottomrule
\end{tabular}
\end{center}

The direction is the predicted one and the mechanism is the predicted
one. \S\ref{sec:microarch} argued that \code{VPDPBUSD} must help the baseline
more than the packed format, because the baseline spends 4 of its 17 vector
operations in the contraction the instruction collapses while the packed kernel
spends 5 of 42. The measurement agrees: $1.24\times$ against $1.15\times$, and
$S$ rises accordingly.

\textbf{The size is modest and the spread is not smaller than the effect.} $S$
rose in six runs of eight, a majority and not a clean separation; the host is a
shared two-vCPU sandbox, and the reviewer's own caveat --- that run-to-run
variation there exceeds the effect being measured --- is adopted rather than
argued away. An \SI{8.8}{\percent} shift in $S$ against $\Sigma$ values ranging
from \num{1.35} to \num{3.03} does not by itself decide
\eqref{eq:condition} either way at any thread count.

What the static model got right, and where it was pessimistic:
\S\ref{sec:microarch} bounded VNNI's cost at \numrange{0}{14}\,\% and this
measurement lands at \SI{8.8}{\percent}, inside that band. The model predicted
$S(\text{vnni}) = \num{2.08}$ on \code{znver4} and \num{2.51} on the two Intel
targets; the measured \num{2.22} sits between them. Given that the model's
negative control failed and its VNNI loops had never been executed, agreement to
this degree is more than was claimed for it.

What remains open is a clean machine --- or a better design, which is
what was done instead. Six of eight is a majority because the eight runs timed
the four kernels in sequence, so the AVX2 and VNNI arms of each ratio came from
different moments on a drifting host. \code{bench\_vnni.c} now interleaves them
for the reason \S\ref{sec:endtoend} interleaves its arms:
31 rounds, all four kernels measured in rapid alternation within each round with
the starting kernel rotated, $S$ formed inside the round, and the result
reported as a median difference with an exact two-sided sign test. That converts
a between-condition comparison on a noisy host into a paired one,
which is the design a shared machine can actually answer; the numbers above
predate it and a dedicated part is no longer the only way to sharpen them. The claim this
subsection supports is narrow: on real VNNI hardware the packed format's
arithmetic disadvantage grows by roughly a tenth, in the direction and for the
reason \S\ref{sec:microarch} predicted, which is a smaller penalty than the
eight-bit storage choice of a system built for those targets would suggest.

\subsection{A second microarchitecture, measured}\label{sec:second}

A reviewer ran the benchmarks on an Intel Xeon (AVX-512 capable, 2 vCPU cloud
sandbox), giving a genuine second data point for everything except VNNI. Their
caveats are adopted rather than paraphrased: with two vCPUs and no pinning, only
the one- and two-thread rows carry information; the four-thread \code{bench\_token}
row and every \code{bench\_threads} row from $T=3$ are oversubscribed and mean
nothing.

\begin{center}
\begin{tabular}{lrr}
\toprule
 & this host (Zen 2) & reviewer (Intel Xeon) \\
\midrule
multiply class, ops/cycle       & \num{1.05}  & \numrange{1.32}{1.44} \\
cheap class, ops/cycle          & \numrange{3.0}{4.0} & \numrange{2.97}{3.04} \\
\code{i2\_s} reference, GMAC/s/core & \num{81.93} & \num{49.86} \\
t5b, GMAC/s/core                & \num{33.74} & \num{23.61} \\
\textbf{$S$}                    & \textbf{\num{2.43}} & \textbf{\num{2.11}} \\
$\Sigma$ at 1 thread            & \num{1.35}  & \num{1.65} \\
$\Sigma$ at 2 threads           & \num{1.63}  & \num{2.23} \\
\bottomrule
\end{tabular}
\end{center}

Two readings, the second load-bearing. First, $S$ does not collapse on a wider
machine: \num{2.11} against \num{2.43}, two runs within \SI{1}{\percent} of each
other. Second, \textbf{this is an out-of-sample test of \S\ref{sec:microarch}'s
model, and it passes.} That model predicted $S = \num{2.098}$ for
\code{icelake-server} and \num{2.317} for \code{sapphirerapids} --- bracketing
the measured \num{2.11} --- from a calibration performed only against this Zen 2
host. A static pipeline model tuned on one microarchitecture predicted a second
it had never seen, to within \SI{1}{\percent}.

At two threads the reviewer's machine gives $S = \num{2.11} < \Sigma =
\num{2.23}$, so \eqref{eq:condition} predicts a marginal win there and none at
one thread. That is a prediction and not a result: no end-to-end run was made on
that host, and the elevated $\Sigma$ comes from the VM's low per-core DRAM
bandwidth (\SI{8.1}{\giga\byte\per\second} single-threaded) rather than from the
microarchitecture.

\paragraph{One model size and one architecture family.} No public ternary
checkpoint larger than 2\,B exists. The synthetic 60-layer graph doubles weight
traffic but not head count, hidden size or vocabulary.

\paragraph{Measurement environment.} The host is shared and run-to-run spread is
comparable to the effect: across eight invocations at ordinary load, one has the
denser code losing. The headline is therefore the five-invocation replication
rather than any single run.

\paragraph{Verification scope.} The round trip is exact on all 210 tensors and
all \num{2084044800} weights, and so --- since this revision --- is the
dot-product comparison: \num{126000} row-and-activation cases across every
\code{i2\_s} tensor, on which the packed kernel matches exact \code{int64}
arithmetic without exception. What remains limited is the read-back from disk,
covering the 210 plus the embedding rather than all 332 tensors; the converter's
own reread covers the dims, type, offset and alignment of all 332.

Two things this does \emph{not} establish, and they are the ones to hold against
it. Nothing here runs the converted model: that the weights survive repacking is
not evidence that a loader reading them emits the same tokens, and it cannot be
until a loader knows type 43. And exhaustiveness over tensors is not
exhaustiveness over inputs --- 200 rows per tensor under three activation
patterns is a large sample of a space that is not finite.

\paragraph{Task benchmarks.} None. The format question does not admit one ---
output is bit-identical --- and where this \emph{model} stands against
others is a different question this paper does not address. We ran a cross-model
perplexity comparison and do not report it as a result, for two reasons, both of
which leave the identity result untouched because both arms are affected
equally.

The first is metadata: the checkpoint's GGUF requires \code{tokenizer.ggml.pre}
and the EOS id to be supplied through \code{-{}-override-kv} because its own
metadata is wrong, so an absolute perplexity from it is at least as likely to
reflect a tokenisation artefact as the model.

The second we found while auditing this work, and it is not a suspicion but a
measurement: the graph \code{llama.cpp} builds for this architecture is not the
model's. \code{build\_ffn} is called with \code{LLM\_FFN\_SILU} in
\code{src/models/bitnet.cpp}, whereas the checkpoint's own \code{config.json}
declares \code{"hidden\_act": "relu2"} and its technical
report~\cite{bitnet2b4t} states that ``instead of the commonly used SwiGLU
activation \dots{} BitNet b1.58 2B4T employs squared ReLU''. The GGUF does not
settle it: it declares \code{general.architecture = bitnet-b1.58} and carries no
activation key at all, so the per-architecture default in the loader is the only
thing choosing, and \code{LLM\_FFN\_RELU\_SQR} exists beside it in
\code{src/llama-graph.h} and is used by five other architectures.

Substituting it is one token of diff, and we ran it. Over the same forty chunks
of WikiText-2:
\begin{center}
\begin{tabular}{lrr}
\toprule
 & \code{i2\_s} & \textbf{t5b} \\
\midrule
\code{LLM\_FFN\_SILU} (as built)   & $\num{96.8243}\pm\num{3.55673}$ & $\num{96.8243}\pm\num{3.55673}$ \\
\code{LLM\_FFN\_RELU\_SQR}         & $\num{16.7482}\pm\num{0.50521}$ & $\num{16.7482}\pm\num{0.50521}$ \\
\bottomrule
\end{tabular}
\end{center}
A factor of \num{5.78}. The \num{96.8243} quoted above is therefore a property
of a mis-built graph and not of this checkpoint, which is the second and larger
of the two reasons it is not reported as a result.

\emph{The identity survives the change, and that was the check that could have
failed.} The two arms agree to six significant figures under \emph{both}
activations --- they have always run the same graph, so a wrong activation
cancels exactly, and a divergence here would have meant something other than the
activation was wrong. Nothing else in this paper is affected: every rate, every
density figure and the bit-identity result are comparisons between two formats
under one loader.

Both of those columns are one tree and two formats neither of which upstream
maintains, which is a weaker position than it first appears, so we repeated the
pair on \code{ggml-org/llama.cpp} at \code{bbdd9f2} using \code{TQ1\_0} ---
\code{llama.cpp}'s own ternary format, and the prior art of \S\ref{sec:contrib}.
Reaching it took two steps, both recorded in
\code{results/activation\_upstream\_master.txt}: \code{i2\_s} is not an upstream
type, and the fork writes an architecture string, \code{bitnet-b1.58}, that
upstream does not register, so the file was retagged to \code{bitnet} after
checking that the fork's subclass overrides only its layer-count-to-label map
and that the two tensor sets agree exactly.
\begin{center}
\begin{tabular}{llrr}
\toprule
tree & format & \code{SILU} & \code{RELU\_SQR} \\
\midrule
this paper's fork   & \code{i2\_s} & \num{96.8243} & \num{16.7482} \\
this paper's fork   & t5b          & \num{96.8243} & \num{16.7482} \\
upstream \code{bbdd9f2} & \code{TQ1\_0} & \num{96.6913} & \num{16.7750} \\
\bottomrule
\end{tabular}
\end{center}
The \code{SILU} arms agree to \SI{0.14}{\percent} and the \code{RELU\_SQR} arms
to \SI{0.16}{\percent} across three quantisations whose only shared component is
the graph, and both upstream arms repeat bit-identically with the load order
alternated. A quantisation defect cannot produce that pattern. We note the
absolute levels are not this checkpoint's under either activation --- every one
of these runs logs a missing pre-tokenizer type --- so what the table supports is
the ratio and not the level.

\paragraph{Engineering.} The blocked matrix--matrix kernel spills 181 times per
loop at its chosen width; narrower widths spill less and measure slower. The
type identifier used (43) is not reserved by upstream, so a stock
\code{llama.cpp} cannot read the file by construction rather than by accident.

\section{Conclusion}

The gap between the \num{1.58} bits a ternary weight carries and the \num{2.00}
it is stored in is recoverable on commodity AVX2 hardware without lookup tables.
Five ternary values fit a byte at exactly \num{1.600} bits; all four radix-3
quotients are extractable from that byte with one independent multiply each,
because the byte is bounded by \num{242} and lies deep inside the exactness range
of 16-bit magic multipliers; and the contraction proceeds on the digits, so no
weight is materialised.

Whether the resulting \SI{25}{\percent} reduction in weight traffic is
profitable is not a property of the code. It is decided
by~\eqref{eq:condition}, a comparison between a core's multiply-port throughput
and its share of memory bandwidth, both measurable in minutes. On the part
measured here the condition holds from four cores upward once the ratio is taken
against the kernel \code{llama.cpp} actually dispatches rather than against its
reference (\S\ref{sec:endtoend}), and the deployed model
is \SI{9}{\percent} smaller and \SIrange{11.9}{13.3}{\percent} faster with
bit-identical output --- the medians over five invocations, prompt and
generation.

The same inequality predicts the result will narrow on a part with VNNI, and
\S\ref{sec:vnnimeas} measures that narrowing at \SI{8.8}{\percent} --- the
direction and mechanism the condition names, at a magnitude the eight-bit
storage choice of a system built for such parts would not have led one to
expect. That is not a caveat appended to a positive result; it is the result,
which is that weight density is a hardware-dependent optimisation with a
computable decision rule rather than a property of the model. What
\S\ref{sec:microarch} does settle is that the rule's input is not parochial:
across four distinct scheduler models spanning two vendors and three
generations, $S$ ranges from \num{2.10} to \num{2.50} and every newer target
sits at or \emph{below} the Zen 2 value it was derived on, by at most
\SI{16}{\percent}. The ratio never worsens on a wider machine, so the rule is
worth applying rather than an accident of the one part that produced it.

\section*{Corrections}

Twenty-nine claims in earlier drafts of this paper were wrong and were corrected
before submission --- among them the novelty of the packing, of the depth-one
extraction arithmetic, of the fused digit-plane contraction and of the
per-tensor scale model, two statements about dead-neuron removal, and the assertion that a
run on VNNI hardware would settle \S\ref{sec:microarch}. Each is listed with
what replaced it and how it was found in \code{CHANGELOG.md} in the repository
below, kept there rather than here so that this document states what holds.

\paragraph{The concessions, canonically.} A concession of prior art is not a
correction in that sense, and the two lists are different lengths, so the
concession list is fixed here and every other site in this paper matches it.
Six constructions this work uses belong to someone else: the base-3 packing,
the $32k+j$ interleave, the depth-one extraction arithmetic, the fused
digit-plane contraction, the per-tensor scale model and the profitability
condition. \S\ref{sec:priorart} states five of them against their owners and
\S\ref{sec:format} the interleave, and the abstract and \S\ref{sec:tq10} name
the same six and no others.

\paragraph{What a $\pm$ means here.}
Three different things, and the document should say which. Around an activation
bound it is the int8 range. After a perplexity it is
\code{llama-perplexity}'s own figure, the standard error of the mean over
chunks: it computes $\sqrt{(\overline{v^{2}}-\bar v^{2})/(n-1)}$ on the
per-chunk negative log-likelihood and scales it by the perplexity. After a
throughput it is \code{llama-bench}'s own figure over its repetitions. None of
them is a confidence interval and none is comparable to another; each is
reproduced verbatim from the tool that printed it, and every pair in this paper
was checked against the evidence file it came from.

\section*{Data and code availability}

The format, the kernels, the benchmarks, the converters and every evidence file
cited here are public under the Apache License~2.0 at
\begin{center}
\url{https://github.com/purpleskulll/bitnet-t5b}
\end{center}
with the layout below. No model weights are distributed and no upstream source
is redistributed verbatim; the integration is applied as a patch to a checkout
the user supplies, and the files that reproduce or adapt an upstream contract
say so in their own headers and in \code{NOTICE}.

\begin{center}
\small
\begin{tabular}{@{}>{\raggedright\arraybackslash}p{0.36\linewidth}>{\raggedright\arraybackslash}p{0.56\linewidth}@{}}
\toprule
claim & file \\
\midrule
code histogram, code~3 absence, entropy & \code{results/tensor\_structure.txt}, \code{results/real\_weights\_check.txt} \\
instruction-port throughput $\pi$ & \code{benchmarks/bench\_ports.c} \\
surplus $\Sigma$ against core count & \code{results/thread\_headroom.txt}, \code{benchmarks/bench\_threads.c} \\
kernel rates, instruction and spill counts & \code{results/weight\_density.txt}, \code{benchmarks/bench\_alu.c} \\
weight traffic of one token & \code{benchmarks/bench\_token.c} \\
end-to-end throughput & \code{results/inference\_t5b.txt} \\
in-situ matmul time, both arms & \code{results/insitu\_symmetric.txt}, \code{results/insitu\_kernel\_rate.txt} \\
matmul share of a token & \code{results/insitu\_kernel\_rate.txt} \\
against \code{TQ1\_0} at equal weights & \code{results/tq1\_0\_comparison.txt}, \code{tools/gguf\_to\_tq1\_0.py} \\
perplexity agreement between the formats & \code{results/perplexity\_t5b.txt} \\
round trip on real weights & \code{results/real\_weights\_check.txt} \\
dead feed-forward neurons & \code{results/dead\_neurons.txt}, \code{tools/dead\_neurons.py} \\
\bottomrule
\end{tabular}
\end{center}

The format and its kernels are \code{src/ternary\_t5b.\{c,h\}} with the
\code{ggml}-facing contract in \code{src/ggml\_t5b\_glue.\{c,h\}} and the suite in
\code{src/test\_t5b.c}; the word-lane variant of \S\ref{sec:lane} is
\code{src/ternary\_t10.\{c,h\}}; the converter is \code{tools/gguf\_to\_t5b.py}.

\paragraph{What the public repository does not contain.} The measurements that
need a model --- \S\ref{sec:results}'s end-to-end table, the perplexity
agreement --- additionally require a \code{llama.cpp} build carrying the type, a
converted GGUF and the checkpoint itself. The integration is supplied as a patch
against a named upstream commit rather than as copied sources, and the
\code{i2\_s} reference kernel used as the baseline is fetched from that commit
by a script rather than redistributed. Everything else --- the packing, the
kernels, the suite, the converters and every evidence file --- builds and runs
from a clone with a compiler and Python alone, and so do five of the eight
benchmarks: \code{bench\_ports}, \code{bench\_probe\_cost}, \code{bench\_gemm},
\code{bench\_vnni} and \code{bench\_vnni\_emu}. The other three ---
\code{bench\_alu}, \code{bench\_threads} and \code{bench\_token} --- compare
against upstream's \code{i2\_s} kernel and therefore additionally need
\code{tools/fetch\_i2s\_reference.sh} to have run; without it \code{make bench}
skips them with a message rather than building most of its targets and exiting
zero.

\begin{thebibliography}{99}
\bibitem{bitnet158} S.~Ma, H.~Wang, L.~Ma, L.~Wang, W.~Wu, P.~Dong, L.~Zhang,
  J.~Xue, F.~Wei. \emph{The Era of 1-bit LLMs: All Large Language Models are in
  1.58 Bits}. arXiv:2402.17764, 2024.
\bibitem{bitnet2b4t} Microsoft. \emph{BitNet b1.58 2B4T Technical Report}.
  arXiv:2504.12285, 2025.
\bibitem{bitnetcpp} J.~Wang, H.~Zhou, T.~Song, S.~Cao, Y.~Xia, T.~Cao, J.~Wei,
  S.~Ma, H.~Wang, F.~Wei. \emph{Bitnet.cpp: Efficient Edge Inference for Ternary
  LLMs}. arXiv:2502.11880, 2025.
\bibitem{onebitinfra} J.~Wang, H.~Zhou, T.~Song, S.~Mao, S.~Ma, H.~Wang, Y.~Xia,
  F.~Wei. \emph{1-bit AI Infra: Part 1.1, Fast and Lossless BitNet b1.58
  Inference on CPUs}. arXiv:2410.16144, 2024.
\bibitem{tmac} J.~Wei, S.~Cao, T.~Cao, L.~Ma, L.~Wang, Y.~Zhang, M.~Yang.
  \emph{T-MAC: CPU Renaissance via Table Lookup for Low-Bit LLM Deployment on
  Edge}. EuroSys, 2025. arXiv:2407.00088.
\bibitem{tq10} compilade. \emph{ggml-quants: ternary packing for TriLMs and
  BitNet b1.58}. llama.cpp pull request \#8151, August 2024.
\bibitem{spectra} T.~Vaidhya, A.~Kaushal, V.~Jain, F.~Couture-Harpin,
  P.~Shishodia, M.~Behbahani, Y.~Nevmyvaka, I.~Rish. \emph{Spectra 1.1: Scaling
  Laws and Efficient Inference for Ternary Language Models}. arXiv:2506.23025,
  2025.

\bibitem{breyerkorf} T.~M.~Breyer, R.~E.~Korf. \emph{1.6-Bit Pattern
  Databases}. Proceedings of the AAAI Conference on Artificial Intelligence
  24(1):39--44, 2010. doi:10.1609/aaai.v24i1.7558.

\bibitem{miraculix} A.~Freudenberg, J.~Vandenplas, M.~Schlather, T.~Pook,
  R.~Evans, J.~ten~Napel. \emph{Accelerated matrix-vector multiplications for
  matrices involving genotype covariates with applications in genomic
  prediction}. Frontiers in Genetics 14:1220408, 2023.
  doi:10.3389/fgene.2023.1220408.

\bibitem{ntru} J.~M.~Schanck et al. \emph{NTRU reference implementation},
  \code{ref-common/pack3.c}. \url{https://github.com/jschanck/ntru}, 2019.
  NIST post-quantum cryptography finalist; vendored through PQClean into
  liboqs.

\bibitem{zukowski} M.~Zukowski, S.~Heman, N.~Nes, P.~Boncz. \emph{Super-Scalar
  RAM-CPU Cache Compression}. Proceedings of the 22nd International Conference
  on Data Engineering (ICDE), 2006. doi:10.1109/ICDE.2006.150.

\bibitem{shannonllm} Y.~Tan, X.~Chen, G.~Alonso, K.-S.~Wong, B.~He.
  \emph{Approaching Shannon Bound with Lossless LLM Weight Compression}.
  arXiv:2606.15789, 2026. To appear, ISCA 2026.

\bibitem{mulapshufb} W.~Mu{\l}a, D.~Lemire. \emph{Faster Base64 Encoding and
  Decoding using AVX2 Instructions}. arXiv:1704.00605, 2017. See also
  W.~Mu{\l}a, \emph{Implementing byte-wise lookup table with PSHUFB}, 2016.

\bibitem{goetschel} S.~G\"otschel, M.~Weiser. \emph{Compression Challenges in
  Large Scale Partial Differential Equation Solvers}. Algorithms 12(9):197,
  2019. Preprint arXiv:1907.00667.

\bibitem{deca} G.~Gerogiannis, S.~Eyerman, E.~Georganas, W.~Heirman,
  J.~Torrellas. \emph{DECA: A Near-Core LLM Decompression Accelerator Grounded
  on a 3D Roofline Model}. MICRO-58, 2025. doi:10.1145/3725843.3756073.
  Preprint arXiv:2505.19349.

\bibitem{aliaga} J.~I.~Aliaga, H.~Anzt, T.~Gr\"utzmacher, E.~S.~Quintana-Ort\'i,
  A.~E.~Tom\'as. \emph{Compression and load balancing for efficient sparse
  matrix-vector product on multicore processors and graphics processing units}.
  Concurrency and Computation: Practice and Experience, 2021.
  doi:10.1002/cpe.6515.

\bibitem{sherry} H.~Huang, D.~Wu, Q.~Hu, G.~Yu, J.~Yang, J.~Zhu, X.~Liu,
  D.~Wu. \emph{Sherry: Hardware-Efficient 1.25-Bit Ternary Quantization via
  Fine-grained Sparsification}. arXiv:2601.07892, 2026.

\bibitem{spectralkv} M.~Hariri, A.~Luo, W.~Chen, T.~Zhang, Q.~Wang, X.~Han,
  V.~Chaudhary. \emph{Quantize What Counts: More for Keys, Less for Values}.
  Findings of the Association for Computational Linguistics: ACL 2026,
  pp.~26393--26420.
\bibitem{litespark} \emph{Litespark Inference on Consumer CPUs: Custom SIMD
  Kernels for Ternary Neural Networks}. arXiv:2605.06485.
\bibitem{deepcompression} S.~Han, H.~Mao, W.~J.~Dally. \emph{Deep Compression:
  Compressing Deep Neural Networks with Pruning, Trained Quantization and
  Huffman Coding}. ICLR, 2016. arXiv:1510.00149.
\bibitem{dfloat11} T.~Zhang et al. \emph{70\% Size, 100\% Accuracy: Lossless LLM
  Compression for Efficient GPU Inference via Dynamic-Length Float}.
  arXiv:2504.11651, 2025.
\bibitem{llamacpp} G.~Gerganov et al. \emph{llama.cpp / ggml}.
\bibitem{wikitext} S.~Merity, C.~Xiong, J.~Bradbury, R.~Socher. \emph{Pointer
  Sentinel Mixture Models}. arXiv:1609.07843, 2016.
\end{thebibliography}

\end{document}
